\documentclass[12pt]{article}

% Compile with LuaLaTeX

\usepackage[margin=1.25in]{geometry}
\usepackage{setspace}
\usepackage{microtype}
\usepackage{titlesec}
\usepackage{parskip}
\usepackage{amsmath}
\usepackage{amssymb}
\usepackage{longtable}
\usepackage{array}

\newcolumntype{A}{>{\begin{otherlanguage}{arabic}\RL}p{3.2cm}<{\end{otherlanguage}}}

\usepackage{hyperref}
\usepackage{tipa}

\onehalfspacing
\usepackage{fontspec}
\usepackage[bidi=basic, main=english]{babel}
\babelprovide[import]{arabic}

\babelfont{rm}{TeX Gyre Termes}
\babelfont[arabic]{rm}[
  Renderer=HarfBuzz,
  Path=fonts/,
  Extension=.ttf,
  UprightFont=*-Regular,
  ItalicFont=*-Italic,
  BoldFont=*-Bold,
  BoldItalicFont=*-BoldItalic,
  Script=Arabic,
  Scale=1.1
]{Amiri}

\newcommand{\AR}[1]{\foreignlanguage{arabic}{#1}}

\hypersetup{
  unicode=true,
  colorlinks=true,
  linkcolor=black,
  urlcolor=blue,
  citecolor=black
}

\title{\textbf{Nigerian Pidgin English:\\
Grammar and Core Vocabulary with Script Representation}}

\author{Flyxion}
\date{\today}

\begin{document}

\maketitle

\begin{abstract}
\noindent
Nigerian Pidgin English is a major contact language spoken by tens of millions of people across Nigeria and neighboring regions of West Africa. Although frequently described in popular discourse as a “pidgin” in the sense of a simplified or auxiliary code, it exhibits systematic phonology, stable morphosyntax, productive aspect marking, and discourse strategies comparable to other creole and post-creole continua. Its grammatical organization cannot be reduced to defective English; rather, it reflects a structured outcome of sustained multilingual contact.

This essay offers a formal introduction to Nigerian Pidgin as a rule-governed linguistic system. It surveys its segmental phonology, pronominal paradigm, tense–aspect system, serial verb constructions, negation patterns, interrogative formation, and high-frequency lexical domains. In addition to descriptive analysis, the study develops parallel representational systems in both Latin phonemic transcription and a constrained Arabic-script transliteration framework. The goal is not conversational fluency but structural orientation: to demonstrate that Nigerian Pidgin possesses internal regularities and distributional constraints sufficient to sustain grammatical description, cross-script encoding, and computational modeling.
\end{abstract}

\newpage
\section{Introduction}

Languages that arise from contact are often mischaracterized as transitional or structurally incomplete. Such descriptions obscure the fact that contact languages stabilize through regularization, constraint, and functional pressure. Once stabilized, they display patterned morphosyntax and systematic phonology no less than historically older literary languages. Nigerian Pidgin is exemplary in this regard.

The purpose of this study is not to rehearse sociolinguistic debates over classification, nor to provide a phrasebook-style inventory of expressions. Instead, it adopts a structural perspective. Nigerian Pidgin is treated here as a linguistic system whose components can be described in terms of phoneme inventories, morphosyntactic categories, and combinatorial rules. Its apparent transparency to speakers of English masks underlying reorganizations of tense–aspect marking, argument structure, and discourse management that are neither accidental nor ad hoc.

The approach taken in the following sections proceeds from foundational description to grammatical architecture. Phonological inventory establishes the sound system. Pronominal forms and aspectual markers reveal core syntactic organization. Serial verb constructions illustrate clause-level compositionality. Negation and interrogative formation demonstrate rule-governed transformations rather than lexical substitution. Vocabulary is presented not as a memorization list but as evidence of semantic structuring within recurrent domains.

By foregrounding structural regularity, this introduction aims to reposition Nigerian Pidgin within linguistic analysis as a coherent grammar rather than a derivative code. Such treatment is essential not only for descriptive accuracy but also for orthographic and computational extension. Accordingly, the study develops parallel Latin phonemic representations and a formally constrained Arabic-script transliteration system. These representational layers do not precede structure; they depend upon it. The primary object of study remains the internal organization that makes systematic cross-script encoding possible.

\section{Historical and Linguistic Background}

Nigerian Pidgin emerged from sustained contact between English speakers and diverse West African language communities beginning in the early colonial period. It developed as a trade and interethnic communication language and gradually stabilized across regions. While English supplies much of the lexical material, the grammar reflects substrate influence from Niger--Congo languages, especially in aspect marking, serial verb constructions, and pronoun systems.

Today Nigerian Pidgin functions as a lingua franca in urban centers such as Lagos, Port Harcourt, Warri, and Benin City. It is used in music, radio broadcasting, informal media, comedy, and everyday interaction. Although it does not yet possess full official status in education, its sociolinguistic vitality is undeniable.

Crucially, Nigerian Pidgin is not ``broken English.'' It is a rule-governed system. Many of its features represent simplification relative to Standard English morphology, but simplification does not imply absence of structure. Instead, complexity is reorganized, particularly in the tense--aspect system.

\section{Phonological Tendencies}

Nigerian Pidgin pronunciation reflects both English input and substrate phonological systems. Certain consonant clusters are simplified. Word-final consonants may be reduced. Vowel quality is more stable and less diphthongized than in many English dialects.

Examples:

\begin{quote}
\textit{school} → \textit{skul} \\
\textit{friend} → \textit{fren} \\
\textit{thing} → \textit{tin}
\end{quote}

The interdental fricatives /θ/ and /ð/ are typically replaced by /t/ or /d/ respectively:

\begin{quote}
\textit{three} → \textit{tree} \\
\textit{that} → \textit{dat}
\end{quote}

Stress patterns are more even, and rhythm often reflects syllable-timed tendencies common in West African languages rather than the stress-timed rhythm typical of British English.

\section{Core Morphosyntactic Structure}

\subsection{Basic Word Order}

Nigerian Pidgin exhibits canonical Subject--Verb--Object (SVO) word order, consistent with English surface structure but lacking the inflectional morphology characteristic of Standard English.

\begin{quote}
\textit{I chop rice.} \\
(I eat rice.)
\end{quote}

Verbs do not inflect for person or number. There is no subject--verb agreement morphology:

\begin{quote}
\textit{I go.} \\
\textit{You go.} \\
\textit{Dem go.}
\end{quote}

The absence of inflection does not reduce grammatical clarity because temporal and aspectual distinctions are encoded analytically through preverbal particles.

\subsection{Pronoun System}

The pronominal system is stable and invariant across most dialectal varieties:

\begin{center}
\begin{tabular}{ll}
\textbf{Singular} & \textbf{Plural} \\
I & We \\
You & Una \\
E (he/she/it) & Dem \\
\end{tabular}
\end{center}

Notable structural features of Nigerian Pidgin further illustrate its grammatical economy and functional distinctiveness. One of the most salient characteristics is the absence of gender distinction in the third person singular. The form *e* serves uniformly for referents that, in Standard English, would require differentiation between *he*, *she*, or *it*. This neutralization of grammatical gender reflects a streamlined pronominal system in which semantic interpretation is resolved contextually rather than morphologically.

Equally significant is the second person plural form *una*, which encodes number overtly in a way that contemporary Standard English no longer does. Whereas English relies on contextual inference or dialectal innovations such as *you all*, Nigerian Pidgin maintains a dedicated plural pronoun, thereby preserving a number distinction that enhances pragmatic clarity in address.

The third person plural form *dem* likewise demonstrates multifunctionality. In addition to serving as an independent plural pronoun, it may operate as a post-nominal plural marker, as in constructions equivalent to “the children them.” This dual role exemplifies a broader tendency within the language toward analytic strategies and flexible category boundaries, where lexical items may extend beyond narrowly delimited grammatical functions.

Example:

\begin{quote}
\textit{Dem dey come.} \\
(They are coming.)
\end{quote}

Plural marking may also appear post-nominally:

\begin{quote}
\textit{Di boy dem.} \\
(The boys.)
\end{quote}

Here \textit{dem} functions as a pluralizer rather than a pronoun.

\subsection{Copula and Existential Constructions}

The copular system differs from English. The verb \textit{be} does not function as an inflected auxiliary. Instead, several strategies appear.

\textbf{Zero copula:}

\begin{quote}
\textit{I happy.} \\
(I am happy.)
\end{quote}

\textbf{Na-cleft structure:}

\begin{quote}
\textit{Na doctor e be.} \\
(It is a doctor.)
\end{quote}

\textit{Na} serves as a focus or identificational marker.

\textbf{Existential construction:}

\begin{quote}
\textit{Get problem.} \\
(There is a problem.)
\end{quote}

or

\begin{quote}
\textit{E get problem.}
\end{quote}

\subsection{Verb Invariance}

Verbs in Nigerian Pidgin do not inflect for tense or agreement. Instead, temporal reference is either contextually inferred or expressed through preverbal markers, discussed in the next section.

Examples:

\begin{quote}
\textit{I chop yesterday.} \\
(I ate yesterday.)

\textit{I go tomorrow.} \\
(I will go tomorrow.)
\end{quote}

Temporal adverbs frequently carry the burden of time specification in the absence of morphological tense.

\section{Tense Aspect and Mood}

Nigerian Pidgin relies primarily on preverbal particles rather than inflectional morphology to express tense, aspect, and modality. These particles occur between the subject and the verb and remain invariant across persons.

\subsection{Progressive Aspect: \textit{dey}}

The particle \textit{dey} marks progressive or ongoing action.

\begin{quote}
\textit{I dey chop.} \\
(I am eating.)

\textit{Dem dey work.} \\
(They are working.)
\end{quote}

The progressive is thus analytic rather than inflectional. The verb itself remains unmodified.

\subsection{Past Marker: \textit{bin}}

The particle \textit{bin} marks anterior or past reference.

\begin{quote}
\textit{I bin go.} \\
(I went / I had gone.)

\textit{Dem bin see am.} \\
(They saw him/her/it.)
\end{quote}

In many contexts, past time may also be inferred from adverbials without overt marking. The use of \textit{bin} tends to emphasize temporal distance or narrative sequencing.

\subsection{Future Marker: \textit{go}}

The particle \textit{go} expresses future reference or intention.

\begin{quote}
\textit{I go come tomorrow.} \\
(I will come tomorrow.)

\textit{E go rain.} \\
(It will rain.)
\end{quote}

Unlike English “will,” \textit{go} does not inflect and does not require auxiliary inversion in questions.

\subsection{Completive Aspect: \textit{don}}

The particle \textit{don} marks completed action, roughly corresponding to the English present perfect or perfective.

\begin{quote}
\textit{I don chop.} \\
(I have eaten / I already ate.)

\textit{Dem don finish.} \\
(They have finished.)
\end{quote}

\textit{Don} emphasizes completion rather than temporal location.

\subsection{Habitual Aspect}

Habitual meaning may be expressed through context or through repeated use of \textit{dey} in certain dialects:

\begin{quote}
\textit{I dey go market every day.} \\
(I go to the market every day.)
\end{quote}

The distinction between progressive and habitual is often context-driven rather than morphologically encoded.

\subsection{Modal Constructions}

Modal meanings are typically expressed lexically rather than through auxiliary inflection.

\textbf{Ability:}

\begin{quote}
\textit{I fit do am.} \\
(I can do it.)
\end{quote}

\textbf{Obligation:}

\begin{quote}
\textit{You for go.} \\
(You should have gone / You ought to go.)
\end{quote}

\textbf{Necessity:}

\begin{quote}
\textit{You must go.} \\
(You must go.)
\end{quote}

The system is thus largely analytic. Grammatical categories are expressed through separate lexical items rather than bound morphology.

\subsection{Particle Ordering}

Particles may combine in structured sequences. For example:

\begin{quote}
\textit{I bin dey work.} \\
(I was working.)

\textit{Dem don dey wait.} \\
(They have started waiting / They are already waiting.)
\end{quote}

The ordering reflects temporal scope, with past or completive markers typically preceding progressive marking.

\section{Negation and Interrogative Structure}

\subsection{Negation}

Negation in Nigerian Pidgin is typically expressed with the preverbal particle \textit{no}.

\begin{quote}
\textit{I no go.} \\
(I did not go / I do not go.)

\textit{Dem no dey work.} \\
(They are not working.)
\end{quote}

Negation precedes aspectual markers when both are present:

\begin{quote}
\textit{I no bin see am.} \\
(I did not see him/her/it.)
\end{quote}

The structure remains analytic and invariant; the verb itself does not change form.

\subsection{Negative Imperatives}

Prohibition is typically formed using \textit{no} or \textit{make no}:

\begin{quote}
\textit{No go there.} \\
(Don't go there.)

\textit{Make you no talk.} \\
(Do not speak.)
\end{quote}

\subsection{Yes--No Questions}

Yes--no questions are generally formed through intonation rather than auxiliary inversion.

\begin{quote}
\textit{You dey come?} \\
(Are you coming?)

\textit{Dem don finish?} \\
(Have they finished?)
\end{quote}

There is no inversion comparable to English “Are you coming?” The surface structure remains declarative, and interrogative force is conveyed prosodically.

\subsection{Wh-Questions}

Interrogative words are typically fronted but do not trigger subject--auxiliary inversion.

\begin{center}
\begin{tabular}{ll}
Wetyn & What \\
Who & Who \\
Where & Where \\
When & When \\
Why & Why \\
How & How \\
\end{tabular}
\end{center}

Examples:

\begin{quote}
\textit{Wetyn you dey do?} \\
(What are you doing?)

\textit{Where dem go?} \\
(Where did they go?)
\end{quote}

Again, the verb remains uninflected.

\subsection{Focus and Emphasis with \textit{na}}

The particle \textit{na} serves as a focus marker and clefting device.

\begin{quote}
\textit{Na me do am.} \\
(It was I who did it.)

\textit{Na today we go start.} \\
(It is today that we will begin.)
\end{quote}

The structure resembles cleft constructions in other creole languages and in certain Niger--Congo substrate languages.

\section{Serial Verb Constructions}

Serial verb constructions are common and reflect substrate influence from West African languages. Multiple verbs may appear in sequence without conjunctions, expressing closely linked actions.

\begin{quote}
\textit{I go buy book.} \\
(I will go and buy a book.)

\textit{Dem carry come.} \\
(They brought it.)

\textit{I take knife cut am.} \\
(I used a knife to cut it.)
\end{quote}

In the final example, \textit{take} functions as an instrumental marker rather than a fully independent verb in the English sense.

Serial constructions allow compact expression of direction, instrumentality, causation, and sequential action without subordination or conjunction.

\section{Core Vocabulary and Lexical Domains}

Although much of Nigerian Pidgin vocabulary derives historically from English, lexical meaning often diverges in scope, usage, or pragmatic force. The lexicon also incorporates substrate and regional elements.

\subsection{Basic Verbs}

A small set of high-frequency verbs carries much of the communicative load:

\begin{center}
\begin{tabular}{ll}
\textit{chop} & eat \\
\textit{go} & go \\
\textit{come} & come \\
\textit{see} & see \\
\textit{hear} & hear \\
\textit{talk} & speak \\
\textit{carry} & carry / bring \\
\textit{take} & take / use \\
\textit{give} & give \\
\textit{make} & make / cause \\
\end{tabular}
\end{center}

Many verbs have broader semantic range than their English equivalents. For example, \textit{chop} may also mean ``consume'' metaphorically:

\begin{quote}
\textit{Dem chop money.} \\
(They embezzled money.)
\end{quote}

\subsection{Common Nouns}

\begin{center}
\begin{tabular}{ll}
\textit{pikin} & child \\
\textit{woman} & woman \\
\textit{man} & man \\
\textit{house} & house \\
\textit{money} & money \\
\textit{wahala} & trouble \\
\textit{palava} & problem / dispute \\
\end{tabular}
\end{center}

\textit{Wahala} and \textit{palava} are widely used to indicate difficulty or complication.

\subsection{Function Words and Particles}

Function words play a critical structural role.

\begin{center}
\begin{tabular}{ll}
\textit{dey} & progressive marker \\
\textit{don} & completive \\
\textit{bin} & past \\
\textit{go} & future \\
\textit{no} & negation \\
\textit{na} & focus marker \\
\textit{wey} & relative marker \\
\end{tabular}
\end{center}

Relative clauses are typically formed using \textit{wey}:

\begin{quote}
\textit{Di man wey come yesterday.} \\
(The man who came yesterday.)
\end{quote}

\subsection{Discourse Markers}

Nigerian Pidgin includes a range of discourse markers used for emphasis, alignment, and pragmatic nuance:

\begin{center}
\begin{tabular}{ll}
\textit{sha} & mild emphasis / concession \\
\textit{o} & emphasis or emotional marking \\
\textit{abeg} & please / request softener \\
\textit{now} & intensifier or focus particle \\
\end{tabular}
\end{center}

Examples:

\begin{quote}
\textit{I go come sha.} \\
(I will come, anyway.)

\textit{Abeg help me.} \\
(Please help me.)

\textit{Na today now.} \\
(It is today indeed.)
\end{quote}

These markers operate pragmatically rather than syntactically, contributing to tone and stance.

\section{Semantic Extension and Pragmatic Usage}

Lexical meaning in Nigerian Pidgin frequently extends beyond its etymological English base. Words may broaden, narrow, or shift in pragmatic value.

For example:

\begin{quote}
\textit{How far?} \\
(Literal: How far?) \\
Meaning: How are things? What is happening?
\end{quote}

\begin{quote}
\textit{You dey?} \\
(Are you present? / Are you available?)
\end{quote}

Meaning often depends on context, shared cultural understanding, and prosodic cues rather than morphological marking.

\section{Grammatical Enrichment}

\subsection{Reduplication}

Reduplication functions as an intensifier, distributive marker, or expressive device. It is common in many West African languages and appears productively in Nigerian Pidgin.

\textbf{Intensification:}

\begin{quote}
\textit{small small} \\
(little by little)

\textit{fine fine} \\
(very fine / extremely beautiful)
\end{quote}

\textbf{Distributive meaning:}

\begin{quote}
\textit{Dem come one one.} \\
(They came one by one.)
\end{quote}

Reduplication adds semantic nuance without inflectional morphology.

\subsection{Comparative and Superlative Structures}

Comparison does not rely on inflectional endings such as “-er” or “-est.” Instead, comparison is expressed analytically.

\textbf{Comparative with \textit{pass}:}

\begin{quote}
\textit{Dis one big pass dat one.} \\
(This one is bigger than that one.)
\end{quote}

\textit{Pass} functions as a comparative marker meaning “more than.”

\textbf{Superlative with \textit{pass all}:}

\begin{quote}
\textit{Na im fine pass all.} \\
(He/she is the most beautiful.)
\end{quote}

Comparison is thus relational rather than morphologically marked.

\subsection{Quantifiers and Plurality}

Plural marking is optional and often contextually determined. The plural marker \textit{dem} may follow a noun.

\begin{quote}
\textit{Book dem dey table.} \\
(The books are on the table.)
\end{quote}

Quantifiers include:

\begin{center}
\begin{tabular}{ll}
\textit{plenty} & many / much \\
\textit{small} & little \\
\textit{all} & all \\
\textit{every} & every \\
\end{tabular}
\end{center}

Examples:

\begin{quote}
\textit{Plenty people come.} \\
(Many people came.)

\textit{Small water remain.} \\
(A little water remains.)
\end{quote}

\subsection{Possession}

Possession is typically expressed through juxtaposition rather than apostrophe marking.

\begin{quote}
\textit{Di man house.} \\
(The man's house.)

\textit{My papa car.} \\
(My father's car.)
\end{quote}

The possessive relation is interpreted structurally rather than morphologically.

\subsection{Relative Clauses}

The relative marker \textit{wey} introduces subordinate descriptive clauses.

\begin{quote}
\textit{Di woman wey dey talk.} \\
(The woman who is speaking.)
\end{quote}

No relative pronoun inflection occurs. The marker remains invariant.

\subsection{Conditional Constructions}

Conditional statements are typically formed with \textit{if} or with serial structures.

\begin{quote}
\textit{If you come, I go happy.} \\
(If you come, I will be happy.)
\end{quote}

Hypothetical or counterfactual nuance may be conveyed through context or through use of \textit{for}.

\begin{quote}
\textit{You for tell me.} \\
(You should have told me.)
\end{quote}

\subsection{Imperatives and Causatives}

Imperatives are generally bare verbs:

\begin{quote}
\textit{Come here.}

\textit{Carry am.}
\end{quote}

Causation frequently uses \textit{make}:

\begin{quote}
\textit{Make I go.} \\
(Let me go.)

\textit{Make dem start.} \\
(Let them begin.)
\end{quote}

Here \textit{make} functions as a causative or permissive marker.

\section{Pragmatics and Politeness Strategies}

Nigerian Pidgin operates within a rich pragmatic ecology in which politeness, hierarchy, solidarity, and emotional stance are frequently conveyed through particles, repetition, tone, and lexical choice rather than through formal grammatical marking.

\subsection{Softening Devices}

The particle \textit{abeg} functions as a politeness marker, mitigating requests and commands.

\begin{quote}
\textit{Abeg give me water.} \\
(Please give me water.)
\end{quote}

It can also appear mid-sentence to soften disagreement:

\begin{quote}
\textit{Abeg, no vex.} \\
(Please, do not be angry.)
\end{quote}

\subsection{Emphatic and Affective Particles}

The sentence-final particle \textit{o} marks emotional emphasis, insistence, or affective engagement.

\begin{quote}
\textit{Come here o.} \\
(Come here, I insist / urgently.)
\end{quote}

The particle \textit{sha} may signal concession or mild contrast:

\begin{quote}
\textit{I go try sha.} \\
(I will try anyway.)
\end{quote}

\subsection{Solidarity and Alignment}

Greetings and discourse openings often establish social alignment.

\begin{quote}
\textit{How far?} \\
(How are things?)

\textit{You dey?} \\
(Are you there? / Are you around?)
\end{quote}

Such expressions function less as literal questions and more as relational markers.

\subsection{Respect and Hierarchy}

Honorific marking is generally lexical rather than morphological. Terms such as \textit{oga} (boss) or \textit{madam} index respect.

\begin{quote}
\textit{Oga, I don finish.} \\
(Boss, I have finished.)
\end{quote}

Politeness strategies rely heavily on intonation, context, and relational expectations.

\section{Code-Switching and English Overlap}

Nigerian Pidgin exists in a dynamic continuum with Standard Nigerian English and with regional languages such as Yoruba, Igbo, and Hausa. Speakers frequently alternate between varieties depending on context, audience, and social positioning.

\subsection{Lexical Mixing}

Code-switching may involve insertion of English technical vocabulary into Pidgin structure.

\begin{quote}
\textit{I go submit di assignment tomorrow.}
\end{quote}

The syntactic frame remains Pidgin while lexical items derive from Standard English.

\subsection{Register Shifting}

In formal contexts, speakers may shift toward Standard English morphology while retaining Pidgin rhythm or particles. In informal contexts, greater reliance on Pidgin structures appears.

This fluidity complicates attempts at strict grammatical boundary drawing. Nigerian Pidgin functions not as an isolated code but as part of a sociolinguistic spectrum.

\section{Regional Variation}

Although broadly intelligible across Nigeria, Nigerian Pidgin exhibits regional phonological and lexical variation.

In southern urban centers such as Warri, Port Harcourt, and Lagos, Pidgin often functions as a primary lingua franca. In northern regions, usage may be more restricted and influenced by Hausa structures.

Lexical items such as \textit{wahala} may carry slightly different pragmatic weight across regions. Pronunciation and intonation patterns vary significantly.

The language is thus internally heterogeneous rather than uniform.

\section{Orthography and Standardization}

Nigerian Pidgin lacks a fully standardized orthography. Spellings vary according to phonetic intuition and influence from English conventions.

For example, the progressive marker may appear as:

\begin{quote}
\textit{dey}, \textit{de}, or \textit{day}
\end{quote}

Similarly, \textit{wetin} may appear as \textit{wetyn}.

Despite this variation, there are emerging efforts toward standardization in media, literature, and education. Broadcast journalism and online platforms increasingly employ consistent spelling conventions.

Orthographic fluidity reflects the primarily oral historical transmission of the language.

\section{Annotated Sample Dialogue}

\begin{quote}
A: \textit{How far? You dey?} \\
B: \textit{I dey. Wetin happen?} \\
A: \textit{I bin call you yesterday but you no pick.} \\
B: \textit{Ah, I bin dey work. I don finish now.} \\
A: \textit{Good. Make we go market small small.} \\
B: \textit{Okay. I go come now.}
\end{quote}

\textbf{Analysis:}

\textit{How far?} functions as greeting rather than literal distance inquiry. \\
\textit{I dey} marks presence or availability. \\
\textit{bin} marks past action. \\
\textit{no} marks negation. \\
\textit{don} marks completion. \\
\textit{make we} forms a hortative construction. \\
\textit{go} marks future intention.

\section{Phonology and an Alternative Orthographic Proposal}

\subsection{Phonological Overview}

Nigerian Pidgin phonology reflects substrate influence from West African languages and superstrate influence from English. The sound system is generally simpler in syllable structure than Standard English, with reduced consonant clusters and strong preference for open syllables in many regional varieties.

Consonant clusters at word-final position are frequently simplified:

\begin{quote}
\textit{ask} $\rightarrow$ \textit{aks} \\
\textit{child} $\rightarrow$ \textit{chail}
\end{quote}

Vowel quality tends toward a reduced inventory compared to English, and diphthongs are often monophthongized depending on speaker background.

Tone is not lexically contrastive in the way it is in Yoruba or Igbo, but intonation carries significant pragmatic meaning, particularly in interrogatives and emphasis.

\subsection{Syllable Structure}

The dominant syllable structure approximates (C)V(C), though final consonants may be weakened or deleted in casual speech. Many borrowings from English undergo phonotactic adaptation to align with West African patterns.

Reduplication and particle-based grammar further reinforce rhythmic and prosodic regularity.

\subsection{The Case for Orthographic Reconsideration}

The current Latin-based orthography inherits English spelling conventions, which introduce inconsistency between sound and symbol. Variants such as \textit{dey}, \textit{de}, and \textit{day} illustrate instability.

An orthography more closely aligned to phonetic realization could reduce ambiguity. One possible direction is transliteration into Arabic script (Ajami-style adaptation), which has historical precedent in West Africa across Hausa and other languages.

Such a system would not imply cultural replacement but phonetic alignment. Arabic script represents consonantal structure efficiently and can be modified with diacritics to capture vowel distinctions.

\subsection{Proposed Arabic-Based Transliteration Framework}

The following is a preliminary mapping based on approximate phonetic correspondence.

Consonants:

\begin{center}
\begin{tabular}{ll}
b & \AR{ب} \\
d & \AR{د} \\
f & \AR{ف} \\
g & \AR{غ} or \AR{گ} (modified) \\
k & \AR{ك} \\
l & \AR{ل} \\
m & \AR{م} \\
n & \AR{ن} \\
p & \AR{پ} (Persian extension) \\
r & \AR{ر} \\
s & \AR{س} \\
t & \AR{ت} \\
w & \AR{و} \\
y & \AR{ي} \\
\end{tabular}
\end{center}

Vowels would be represented through short vowel diacritics and long vowel letters:

\begin{center}
\begin{tabular}{ll}
a & \AR{َ} \\
i & \AR{ِ} \\
u & \AR{ُ} \\
long a & \AR{ا} \\
long i & \AR{ي} \\
long u & \AR{و} \\
\end{tabular}
\end{center}

\subsection{Illustrative Examples}

\textit{I dey} could be rendered:

\begin{quote}
\AR{ا دي}
\end{quote}

\textit{Wetin happen?}

\begin{quote}
\AR{وَتِن هَپَن}
\end{quote}

\textit{I no go}

\begin{quote}
\AR{ا نو غو}
\end{quote}

This transliteration prioritizes phonetic transparency over etymological fidelity.

\subsection{Advantages and Challenges}

The proposed orthographic framework offers several theoretical and practical advantages. Most notably, it enables a closer alignment between phonological structure and graphic representation, thereby reducing the distance between sound and symbol that often complicates the transcription of Nigerian Pidgin. In addition, by situating the system within the broader intellectual and historical continuum of Ajami literacy traditions across West Africa, the proposal opens a pathway for continuity rather than rupture. Such alignment may facilitate cross-linguistic familiarity among communities already accustomed to Arabic-derived scripts, while also supporting the development of locally grounded literacy practices.

At the same time, significant challenges remain. Nigerian Pidgin exhibits substantial dialectal variation, particularly in its vowel realizations, and any fixed orthographic mapping risks privileging one regional norm over others. The absence of a standardized tonal marking system further complicates matters, especially in contexts where prosodic distinctions may carry semantic weight. Beyond linguistic considerations, script reform inevitably intersects with sociopolitical sensitivities. Orthography is rarely neutral; it encodes histories of power, education, religion, and identity. Any attempt to formalize or modify script usage must therefore proceed with attentiveness to community reception and cultural meaning. Finally, technical questions of encoding, font support, keyboard layouts, and digital interoperability pose nontrivial barriers to widespread adoption.

For these reasons, the present proposal is best understood as exploratory rather than prescriptive. Its aim is not cultural substitution, nor the displacement of existing Latin-based conventions, but the pursuit of phonological coherence and structural clarity within an expanded graphic repertoire.

\subsection{Theoretical Implication}

Orthography shapes perception of linguistic structure. A phonetic Arabic-based script would foreground syllabic rhythm and particle-based grammar rather than English lexical ancestry. It would visually decouple Nigerian Pidgin from Standard English, potentially reinforcing its autonomy as a linguistic system.

\section{Toward a Phonetic Ajami Orthography for Nigerian Pidgin}

\subsection{Design Principles}

The proposed Arabic-based transliteration system is not etymological but phonetic. Its primary objective is one-to-one mapping between sound and symbol, thereby minimizing the orthographic opacity inherited from English.

Three structural principles guide the system:

First, consonants are mapped according to phonetic similarity rather than English spelling convention.

Second, vowels are represented consistently through diacritics and long-vowel letters, prioritizing phonemic transparency.

Third, grammatical particles are preserved visually as independent units, reinforcing the analytic structure of Nigerian Pidgin.

\subsection{Consonant Mapping}

The consonant inventory largely overlaps with Arabic but requires extended letters for sounds not native to Classical Arabic. Persian extensions such as \AR{پ} for /p/ are employed where necessary to preserve phonemic contrast.

For example:

\begin{center}
\begin{tabular}{ll}
b & \AR{ب} \\
d & \AR{د} \\
f & \AR{ف} \\
g & \AR{غ} / \AR{گ} \\
k & \AR{ك} \\
l & \AR{ل} \\
m & \AR{م} \\
n & \AR{ن} \\
p & \AR{پ} \\
r & \AR{ر} \\
s & \AR{س} \\
t & \AR{ت} \\
w & \AR{و} \\
y & \AR{ي} \\
h & \AR{ه} \\
\end{tabular}
\end{center}

Affricates such as /ch/ and /j/ may be represented using adapted forms, for example \AR{تش} and \AR{ج}, respectively.

\subsection{Vowel Representation}

Short vowels are indicated through diacritics:

\[
a = \AR{َ} \qquad
i = \AR{ِ} \qquad
u = \AR{ُ}
\]

Long vowels are represented through \AR{ا}, \AR{ي}, and \AR{و}.

For example:

\begin{quote}
\AR{آيْ دي} \\
(I dey)
\end{quote}

Here \AR{آيْ} captures the diphthongal realization of “I,” while \AR{دي} captures the progressive particle /deɪ/.

\subsection{Grammatical Particles in Script}

A significant structural advantage of this orthography is the visual stabilization of analytic particles.

\begin{quote}
\AR{نَا سُوْ} \\
(Na so)

\AR{دِمْ دُونْ كَمْ} \\
(Dem don come)
\end{quote}

Particles such as \AR{نَا} (na), \AR{دُونْ} (don), and \AR{دِي} (dey) become visually discrete units, reinforcing their grammatical independence.

This representation visually encodes the analytic structure of tense and aspect marking more transparently than English-based spelling.

\subsection{Phonotactic Adaptation}

The orthographic system resolves cluster reduction by reflecting actual spoken realization rather than English orthography.

For example:

\begin{quote}
\AR{وِيتِنْ} \\
(Wetin)
\end{quote}

This avoids misleading English silent letters and aligns orthography more closely with speech rhythm.

\subsection{Semantic Autonomy Through Script}

By removing English orthographic anchoring, the script visually decouples Nigerian Pidgin from Standard English. This has structural consequences.

The language appears as a system in its own right rather than a deviation from English. Particles gain structural salience. Prosodic rhythm becomes more apparent.

Orthographic reform therefore does not merely change script; it reorganizes perception of grammatical structure.

\subsection{Limitations and Technical Questions}

Dialectal vowel variation presents challenges. The distinction between /e/ and /ɛ/ or between /o/ and /ɔ/ may require additional diacritic marking if phonemic contrast is deemed significant.

Digital encoding requires standardization of extended letters such as گ and پ across platforms.

Sociolinguistic reception would likely vary across regions and communities.

\subsection{Algorithmic Transliteration}

The system is suitable for deterministic transliteration because it operates primarily on phoneme-to-grapheme mapping rather than etymological mapping.

A bidirectional algorithm would require:

A phonemic tokenizer for Latin Pidgin input.

A rule-based consonant mapping.

A vowel diacritic insertion rule sensitive to syllable structure.

Optional tone marking if future refinement demands it.

The resulting system is structurally simpler than English orthography because it does not encode historical spelling residue.

\section{Phoneme Inventory and IPA Mapping}

This section proposes a working phoneme inventory for Nigerian Pidgin suitable for orthographic design and then specifies an IPA-to-script mapping consistent with the transliteration practice illustrated in the corpus above. The goal is not to legislate a single ``true'' pronunciation for all regions, but to provide a coherent baseline inventory that supports deterministic transliteration.

\subsection{Working Phoneme Inventory}

Nigerian Pidgin exhibits regional variation, but a practical orthography can be built around a core consonant system and a modest vowel system. The inventory below is intended as an operational compromise: fine-grained enough to support stable transcription, coarse-grained enough to remain dialect-robust.

\paragraph{Consonants}

\[
/p,\; b,\; t,\; d,\; k,\; g,\;
m,\; n,\; \eta,\;
f,\; v,\; s,\; z,\; h,\;
w,\; j,\;
l,\; r,\;
tʃ,\; dʒ/
\]

Interdental fricatives /\texttheta/ and /\dh/ are excluded, since they are not phonemic in most Nigerian Pidgin varieties and are typically realized as alveolar stops, as reflected in forms such as \textit{dat} and \textit{tree}.

\paragraph{Vowels}

A five-vowel system provides stable coverage:

\[
/i,\; e,\; a,\; o,\; u/
\]

A seven-vowel refinement with /ɛ/ and /ɔ/ may be optionally encoded, but orthographic stability does not require it. Diphthongs are treated as vowel sequences, with /ai/ especially frequent.

\subsubsection{Latin Baseline (Phonemic)}

Before transliteration into Arabic script, phonemic forms are normalized into a narrow Latin baseline that avoids English spelling conventions. This baseline stabilizes input and prevents interference from English orthography.

\begin{center}
\begin{tabular}{lll}
IPA & Latin & Example \\
\hline
/p/ & p & pikin \\
/b/ & b & boy \\
/t/ & t & tin \\
/d/ & d & dey \\
/k/ & k & ko \\
/g/ & g & go \\
/tʃ/ & ch & chop \\
/dʒ/ & j & enjoy \\
/f/ & f & fit \\
/v/ & v & vex \\
/s/ & s & so \\
/z/ & z & zaa (dialectal) \\
/h/ & h & how \\
/m/ & m & make \\
/n/ & n & no \\
/ŋ/ & ng & sing \\
/l/ & l & lie \\
/r/ & r & far \\
/w/ & w & wetin \\
/j/ & y & you \\
\end{tabular}
\end{center}

This level is purely phonemic. English spellings such as \textit{know}, \textit{three}, or \textit{that} are excluded in favor of \textit{no}, \textit{tri}, and \textit{dat}.
\subsubsection{Arabic Script Mapping}

Arabic transliteration is derived mechanically from the Latin baseline. Only standard Arabic letters are used.

\paragraph{Consonants}

\begin{center}
\begin{tabular}{lll}
IPA & Script & Note \\
\hline
/p/ & \AR{ب} & approximated as /b/ \\
/b/ & \AR{ب} & direct \\
/t/ & \AR{ت} & direct \\
/d/ & \AR{د} & direct \\
/k/ & \AR{ك} & direct \\
/g/ & \AR{ج} & phonetic approximation \\
/f/ & \AR{ف} & direct \\
/v/ & \AR{ف} & merged with /f/ \\
/s/ & \AR{س} & direct \\
/z/ & \AR{ز} & direct \\
/h/ & \AR{ه} & direct \\
/m/ & \AR{م} & direct \\
/n/ & \AR{ن} & direct \\
/ŋ/ & \AR{نْج} & cluster representation \\
/l/ & \AR{ل} & direct \\
/r/ & \AR{ر} & direct \\
/w/ & \AR{و} & direct \\
/j/ & \AR{ي} & direct \\
/tʃ/ & \AR{تش} & digraph \\
/dʒ/ & \AR{ج} & direct \\
\end{tabular}
\end{center}

No additional Persian letters are introduced. All approximations are handled internally through clustering or phonemic merger.

\paragraph{Vowels}

Short vowels are represented by diacritics when phonological clarity requires explicit marking. The vowel /a/ is written with \AR{َ} (fatha), /i/ and /e/ with \AR{ِ} (kasra), and /u/ and /o/ with \AR{ُ} (damma). 

In unstressed environments or highly predictable contexts, short vowels may be omitted in fluent writing; however, full marking ensures deterministic back-mapping to the Latin baseline and eliminates ambiguity in pedagogical contexts.

Long vowels are represented by the corresponding matres lectionis. The vowel /aː/ is written with \AR{ا}, /iː/ with \AR{ي}, and /uː/ with \AR{و}. Where contrastive vowel length is not phonemic, these long-vowel letters function primarily as quality markers rather than as duration markers.

Diphthongs are written as ordered sequences of vowel letters. The diphthong /ai/ is represented as \AR{آي}, while /au/ is written as \AR{آو}. In environments where a glottal onset is phonetically salient, initial \AR{آ} provides a stable orthographic anchor.

\subsubsection{Worked Derivations}

The full transliteration pipeline proceeds from surface form to phonological representation, from phonology to Latin baseline, and finally from the Latin baseline to Arabic script.

\paragraph{Progressive}

The form \textit{I dey} is analyzed phonologically as /ai de/, normalized to the Latin baseline \textit{ai de}, and rendered orthographically as \AR{آيْ دِي}.

\paragraph{Negation}

The expression \textit{I no know} corresponds phonologically to /ai no no/, normalizes to \textit{ai no no}, and is written as \AR{آيْ نُو نُو}.

\paragraph{Wh-question}

The form \textit{Wetin dey happen} is analyzed as /wetin de apen/, normalized to \textit{wetin de apen}, and rendered as \AR{وِيتِنْ دِي أَبِن}.

Consonant clusters are either resolved through vowel insertion consistent with attested pronunciation or indicated with sukūn, as in \AR{ْ}, where pedagogically useful to preserve segmental transparency.

\subsubsection{Functional Stability}

High-frequency grammatical particles maintain fixed orthographic forms. The progressive marker \textit{dey} is written as \AR{دِي}, the completive marker \textit{don} as \AR{دُونْ}, the copular or focus particle \textit{na} as \AR{نَا}, and the negator \textit{no} as \AR{نُو}. 

Stability at this level reinforces morphosyntactic transparency and ensures that analytic tense and aspect marking remains visually distinct.

Stability at this level reinforces morphosyntactic transparency and reduces ambiguity in analytic tense and aspect marking.

\subsubsection{Orthographic Principle}

The system is strictly phonological. Each grapheme corresponds to a surface phoneme in the working inventory. English spelling conventions are excluded from derivation. Dialectal variation can be layered through optional diacritics without altering the core mapping.

The result is a closed, internally consistent transliteration framework: inventory → Latin phonemic normalization → Arabic graphemic realization. Each stage precedes the next.

\subsection{Deterministic Transliteration Rules}

This subsection specifies a deterministic procedure for transliterating a restricted Latin baseline into an Arabic-only orthography. The procedure is designed to be implementable as a finite sequence of rewrite rules. It assumes that the Latin input is already normalized to a phonetic baseline rather than English spelling.

\subsubsection{Input Normalization}

Let the input string be a sequence of tokens separated by spaces. Each token is assumed to be lower-case Latin characters drawn from the set
\[
\{a,e,i,o,u,b,d,f,g,h,j,k,l,m,n,r,s,t,w,y\}
\]
together with the digraphs \texttt{sh} and \texttt{ch} and apostrophe-free orthography. If the source uses English-residue spellings, these must be normalized first.

A minimal normalization map that removes common English residue is:

\[
\texttt{th} \mapsto \texttt{t},\quad
\texttt{ph} \mapsto \texttt{f},\quad
\texttt{ck} \mapsto \texttt{k},\quad
\texttt{qu} \mapsto \texttt{kw}.
\]

This stage is not a model of pronunciation; it is a structural stabilizer for transliteration.

\subsubsection{Tokenization}

Let the token sequence be $T_1,\dots,T_N$. Transliteration applies tokenwise. Punctuation may be retained and copied directly after the corresponding token.

\subsubsection{Consonant Mapping}

Within each token, consonantal segments are mapped deterministically by the following substitution table.

\begin{center}
\begin{tabular}{lll}
Latin & Arabic & Comment \\
\hline
b & \AR{ب} & \\
d & \AR{د} & \\
f & \AR{ف} & \\
g & \AR{ج} & Arabic-only approximation \\
h & \AR{ه} & \\
j & \AR{ج} & shares glyph with g; resolved by context \\
k & \AR{ك} & \\
l & \AR{ل} & \\
m & \AR{م} & \\
n & \AR{ن} & \\
r & \AR{ر} & \\
s & \AR{س} & \\
t & \AR{ت} & \\
w & \AR{و} & consonantal /w/ \\
y & \AR{ي} & consonantal /y/ \\
sh & \AR{ش} & \\
ch & \AR{تش} & Arabic-only cluster encoding \\
\end{tabular}
\end{center}

The map is applied greedily left-to-right with longest-match priority, so \texttt{sh} and \texttt{ch} are rewritten before single-letter consonants.

\subsubsection{Vowel Strategy}

Two regimes are possible. The first is a readability-first regime that writes vowels as letters where possible. The second is a precision-first regime that uses diacritics systematically. The transliteration system in this essay adopts the readability-first regime as default, with diacritics treated as optional refinement.

The readability-first vowel map is:

\begin{center}
\begin{tabular}{lll}
Latin & Arabic & Comment \\
\hline
a & \AR{ا} & default long \\
i & \AR{ي} & default long \\
u & \AR{و} & default long \\
e & \AR{ي} & nearest-grapheme approximation \\
o & \AR{و} & nearest-grapheme approximation \\
\end{tabular}
\end{center}

This strategy produces a stable orthography without requiring diacritic entry. It intentionally sacrifices fine vowel quality distinctions, which in Nigerian Pidgin are frequently predictable from lexical and syntactic context.

An optional refinement permits systematic short-vowel marking in pedagogical mode. Under this refinement, the vowel /a/ is marked with \AR{َ}, the vowels /i/ and /e/ with \AR{ِ}, and the vowels /u/ and /o/ with \AR{ُ}. Such marking may be employed when the writer wishes to avoid ambiguity or to ensure fully deterministic phonological recovery.

\subsubsection{Special Lexical Tokens}

Certain high-frequency function words are best treated lexically in order to preserve recognizable forms and avoid overgeneration. Define a small lexical override dictionary $\mathcal{D}$, applied prior to general rules:

\begin{center}
\begin{tabular}{ll}
Latin & Arabic \\
\hline
i & \AR{آي} \\
na & \AR{نا} \\
no & \AR{نو} \\
dey & \AR{دي} \\
don & \AR{دون} \\
dem & \AR{ديم} \\
wetin & \AR{ويتين} \\
abeg & \AR{ابج} \\
\end{tabular}
\end{center}

This procedure ensures that analytic particles remain visually discrete, reinforcing their grammatical independence and supporting the pedagogical function of the script.

\subsubsection{Ambiguity Management}

Because Arabic-only constraints collapse certain contrasts, two systematic ambiguities arise.

The first is the collapse of /g/ and /j/ under \AR{ج}. Resolution is lexical: in practice, Nigerian Pidgin has relatively few minimal pairs dependent on this contrast, and the intended reading is recoverable from context. If disambiguation is required, one may optionally reserve \AR{غ} for /g/ and \AR{ج} for /j/ as an internal convention, while still remaining within Arabic letters.

The second is the collapse of /p/ into \AR{ب} if the Latin baseline includes \texttt{p}. Under Arabic-only constraints, \texttt{p} is rewritten as \AR{ب}. This is historically normal in many Ajami systems. If one wishes to preserve /p/ without non-Arabic letters, the only option is a diacritic convention on \AR{ب}, but this is best treated as a later refinement rather than a core requirement.

\subsubsection{Algorithmic Order}

The transliteration pipeline is therefore:

First, apply lexical overrides $\mathcal{D}$ tokenwise. Second, apply normalization to remove English residue patterns. Third, apply greedy consonant mapping with longest-match priority. Fourth, apply the vowel strategy, either readability-first or pedagogical-diacritic mode. Finally, reattach punctuation and preserve spacing.

This order ensures that high-frequency particles stabilize early, consonantal structure is preserved deterministically, and vowels are added in a predictable manner.

\subsubsection{Worked Derivations}

\paragraph{Example 1: \texttt{i dey}}

Application of the lexical override dictionary yields \texttt{i} → \AR{آي} and \texttt{dey} → \AR{دي}. The resulting form is therefore \AR{آي دي}.

\paragraph{Example 2: \texttt{wetin dey happen}}

Lexical override first yields \texttt{wetin} → \AR{ويتين} and \texttt{dey} → \AR{دي}. The remaining token \texttt{happen} is processed by consonant mapping, where h corresponds to \AR{ه}, p to \AR{ب}, and n to \AR{ن}. Under the readability-first vowel strategy, the stable orthographic output is \AR{هابن}. The complete expression is therefore rendered as \AR{ويتين دي هابن؟}.

\paragraph{Example 3: \texttt{dem don come}}

Lexical override yields \texttt{dem} → \AR{ديم} and \texttt{don} → \AR{دون}. The remaining token \texttt{come} undergoes consonant mapping, where c corresponds to \AR{ك} and m to \AR{م}. Vowel mapping inserts \AR{و} for /o/ and \AR{ي} for /e/, producing the intermediate form \AR{كومي}. A readability adjustment may optionally omit the final \AR{ي} to reflect common pronunciation, yielding \AR{كوم}. The full expression is thus written as \AR{ديم دون كوم}.

The derivations illustrate that the system is deterministic at the rule level while still permitting a controlled layer of phonological smoothing for high-frequency words.

\section{Comparison with Hausa Ajami}

Any proposal for an Arabic-based orthography for Nigerian Pidgin necessarily invites comparison with Hausa Ajami, the long-standing Arabic-script tradition used to write Hausa across northern Nigeria and the wider Sahel. The comparison is instructive because Hausa Ajami represents one of the most historically successful adaptations of Arabic script to a non-Arabic phonological system in West Africa.

\subsection{Historical Orientation}

Hausa Ajami emerged within Islamic scholarly networks and was used for poetry, legal documents, correspondence, and religious instruction. It developed organically over centuries, rather than through centralized planning. Orthographic conventions stabilized gradually through repeated manuscript transmission.

In contrast, the proposed Pidgin Ajami system is explicitly planned. It begins with a phonological inventory and a deterministic transliteration algorithm. The Hausa system evolved bottom-up; the Pidgin system is top-down.

\subsection{Phonological Adaptation}

Hausa contains phonemic contrasts not present in Arabic, including ejectives and implosives, as well as distinct /p/ and /g/ sounds. Hausa Ajami resolved these gaps in several ways.

First, some traditions introduced modified letters through dotting conventions.

Second, some contrasts were collapsed orthographically and disambiguated contextually.

Third, vowel length distinctions were often marked more carefully than in unvowelled Arabic.

The proposed Pidgin Ajami system aligns most closely with the second strategy. Under Arabic-only constraints, /p/ merges with \AR{ب} and /g/ may be represented by \AR{ج} or \AR{غ}. As in Hausa Ajami, phonological recovery is expected to occur through lexical familiarity rather than perfect graphemic fidelity.

The structural similarity lies in functional tolerance for controlled ambiguity.

\subsection{Vowel Representation}

Hausa distinguishes long and short vowels phonemically. In many Ajami manuscripts, long vowels are written explicitly with \AR{ا}, \AR{ي}, and \AR{و}, while short vowels may be omitted or marked with diacritics depending on genre.

Nigerian Pidgin does not rely heavily on vowel length contrasts for lexical differentiation. The proposed Pidgin Ajami therefore adopts a readability-first vowel strategy, writing core vowels explicitly while permitting optional diacritics for pedagogical precision.

Both systems therefore prioritize functional literacy over phonetic exhaustiveness.

\subsection{Morphological Transparency}

Hausa is morphologically richer than Nigerian Pidgin. It contains noun class remnants, plural patterns, and verbal derivations. Hausa Ajami thus encodes morphological alternations within stems.

Nigerian Pidgin, by contrast, is analytically structured. Tense, aspect, and modality are particle-based rather than inflectional. This produces a different orthographic opportunity.

In Pidgin Ajami, particles such as \AR{نا} (na), \AR{دي} (dey), \AR{دون} (don), and \AR{نو} (no) become visually stable units. Because these particles are invariant, the script reinforces analytic grammar rather than morphological alternation.

Hausa Ajami must accommodate internal stem alternation. Pidgin Ajami highlights syntactic particles.

\subsection{Sociolinguistic Position}

Hausa Ajami historically functioned alongside Arabic literacy and later alongside Latin-script Hausa. It did not replace Latin orthography but coexisted with it.

A Pidgin Ajami system would likely occupy a similar sociolinguistic niche: supplemental rather than replacement orthography. It could function in artistic, pedagogical, or cultural contexts without displacing established Latin usage.

The comparison therefore suggests that script plurality is not structurally destabilizing. West African linguistic ecology already supports dual-script traditions.

\subsection{Structural Divergence}

Despite parallels, there is a critical difference.

Hausa is an indigenous Niger–Congo language with a long-standing literary identity. Ajami writing reinforced an existing linguistic core.

Nigerian Pidgin is a contact language with English lexical ancestry and multi-ethnic substrate influence. Writing it in Arabic script visually decouples it from English orthography. This is not merely phonetic adaptation; it is perceptual reclassification.

Hausa Ajami preserves Hausa as Hausa. Pidgin Ajami repositions Pidgin as structurally autonomous.

\subsection{Implications for Standardization}

The Hausa case demonstrates that:

Orthographic conventions stabilize through repeated communal usage rather than abstract design alone.

Phoneme–grapheme mismatch does not prevent literacy if functional load remains manageable.

Script change does not erase Latin literacy but can coexist with it.

For Nigerian Pidgin, the lesson is not that a perfect mapping is required. Rather, consistency and repeatability are more important than exhaustiveness. A deterministic transliteration framework, once adopted in sustained textual practice, would likely converge toward stable norms just as Hausa Ajami did historically.

\section{Deterministic Reverse Transliteration into English}

This section formalizes the reverse mapping from the Arabic-based orthography into a normalized English phonetic baseline. The goal is structural symmetry: transliteration must be reversible at the phonemic level even if etymological English spelling is not perfectly reconstructed.

\subsection{Scope}

The reverse system maps Arabic grapheme sequences into a normalized English phonetic representation approximating contemporary pronunciation rather than historical spelling.

For example, \AR{كَابِتَالِسْتْ} is rendered as \textit{kapitalist}, rather than the etymologically conventional spelling ``capitalist.'' The objective is phonological transparency rather than historical orthographic fidelity.

not necessarily the etymologically correct ``capitalist.''

\subsection{Preprocessing}

Let the input be a Unicode Arabic string. Processing proceeds tokenwise, separated by whitespace and punctuation.

Arabic diacritics (\AR{َ}, \AR{ِ}, \AR{ُ}, \AR{ْ}) are first stripped unless pedagogical vowel precision is required. This ensures consistent mapping independent of diacritic usage.

\subsection{Consonant Mapping}

The base consonant reversal map is:

\begin{center}
\begin{tabular}{lll}
Arabic & Latin & Notes \\
\hline
\AR{ب} & b & \\
\AR{ت} & t & \\
\AR{د} & d & \\
\AR{ك} & k & \\
\AR{ج} & g or j (contextual) & \\
\AR{ف} & f & \\
\AR{س} & s & \\
\AR{ز} & z & \\
\AR{ه} & h & \\
\AR{م} & m & \\
\AR{ن} & n & \\
\AR{ر} & r & \\
\AR{ل} & l & \\
\AR{و} & w (consonant) & \\
\AR{ي} & y (consonant) & \\
\AR{ش} & sh & \\
\AR{تش} & ch & \\
\AR{غ} & g (if used as g marker) & \\
\end{tabular}
\end{center}

Context rule for \AR{ج}:

If followed by front vowel marker (\AR{ي} or kasra \AR{ِ}), render as j.

Otherwise default to g.

This resolves most ambiguity without dictionary lookup.

\subsection{Vowel Mapping}

Long vowels are interpreted according to the following correspondences. The letter \AR{ا} maps to the vowel \textit{a}, \AR{ي} to \textit{i}, and \AR{و} to \textit{u}. When vowel letters occur between consonants, they are interpreted as full vowels rather than glides, unless preceded by another vowel that licenses glide interpretation.

For example, the form \AR{رُولِنْزْ} is rendered as \textit{rulins} under direct phonetic mapping. An optional smoothing rule permits the sequence \textit{u} to be rendered as \textit{oo} when a more conventional English-like representation is desired, yielding forms such as \textit{rulins} or \textit{rulings}, depending on orthographic preference.

\subsection{Cluster Reconstruction}

Because Arabic orthography may compress or partially neutralize consonant clusters, a subsequent smoothing stage may restore likely English cluster patterns. For instance, the form \AR{سِمْبْتُمْزْ} yields the direct transliteration \textit{simbtums}. A heuristic cluster-repair rule may then apply, such as the transformation \textit{mbt} → \textit{mpt}, producing \textit{simptoms}. A final normalization step may restore conventional spelling, yielding \textit{symptoms}. 

This reconstruction stage is explicitly heuristic and optional. It is not required for phonetic transparency but may be applied when conventional orthography is the desired output.

\subsection{Lexical Overrides}

As with forward transliteration, high-frequency function words may be stabilized through lexical overrides prior to general rule application. For example, \AR{ذَا} may be rendered as \textit{the}, \AR{إِنْ} as \textit{in}, \AR{أُفْ} as \textit{of}, \AR{أَنْدْ} as \textit{and}, and \AR{إِتْ} as \textit{it}. 

These overrides permit restoration of conventional English articles, prepositions, and conjunctions, thereby increasing readability while preserving algorithmic determinism.

\subsection{Worked Example}

Consider the input \AR{كَابِتَالِسْتْ يُتُوبِيَانِيزْمْ}. 

The first processing stage optionally strips diacritics, yielding \AR{كابتالست يتوبيانيزم}. Subsequent stages apply consonant and vowel mapping, followed by optional smoothing and lexical normalization, to produce a phonetic English baseline approximating contemporary pronunciation rather than historical spelling.

Step 2: Consonant map:

k a b t a l s t y t u b y a n i z m

Step 3: Vowel restoration:

kapitalist utopianizm

Step 4: English smoothing:

capitalist utopianism

Final output:

Capitalist Utopianism

\section*{Conclusion}

This study has proposed a structurally grounded introduction to Nigerian Pidgin through three interlocking lenses: grammatical description, phonological analysis, and orthographic engineering. Rather than treating the language as an informal derivative of English, the discussion has approached it as a system with its own internal regularities, aspectual structure, pronominal economy, and discourse pragmatics. The grammatical patterns examined demonstrate that Nigerian Pidgin is not a simplified variant of English but a rule-governed language whose apparent transparency conceals a coherent underlying architecture.

The orthographic proposal extends this structural orientation into the domain of script. By restricting the writing system to the Classical Arabic grapheme inventory and introducing minimal, explicitly defined digraph conventions, the system preserves typographic economy while maintaining phonological intelligibility. The result is neither an imitation of Arabic morphology nor an attempt to reproduce English spelling conventions. It is instead a phonologically motivated encoding scheme that reflects acoustic structure rather than etymology.

Crucially, the transliteration architecture has been formulated in explicitly compositional terms. The movement from Arabic graphemes to phonetic English and from phonetic English to standardized orthography demonstrates that transliteration is not a single act but a layered process. Each layer performs a distinct function: acoustic reconstruction, lexical normalization, and literary smoothing. By separating these stages formally, the system avoids conflating sound with spelling and representation with convention.

The broader significance of this framework lies in its methodological posture. A writing system need not be inherited to be legitimate; it may be engineered, provided its design principles are explicit. When phoneme inventories are clearly specified, grapheme correspondences are deterministically defined, and reading rules are computationally describable, orthography becomes a domain of formal design rather than historical accident.

Nigerian Pidgin, long transmitted primarily through speech and Latin-based informal writing, proves adaptable to such formalization. The exercise confirms that the language possesses sufficient structural stability to sustain cross-script representation without loss of grammatical identity. In doing so, it demonstrates that linguistic description, orthographic construction, and computational modeling need not be separate enterprises. They may converge in a single system grounded in phonology, constrained by script economy, and articulated through formal mapping.

What emerges is not merely a pedagogical experiment but a demonstration of principle: that structure precedes symbol, and that symbol, when responsibly engineered, can faithfully reflect the structure it encodes.

\newpage
\section*{Appendices}

\appendix
\section{Appendix A: Phoneme Inventory of Nigerian Pidgin}

\subsection{Methodological Note}

The phoneme inventory presented here reflects widely attested urban Nigerian Pidgin varieties. Variation exists regionally, particularly under substrate influence from Yoruba, Igbo, and Hausa; the inventory below represents a conservative shared core sufficient for orthographic modeling.

\subsection{Consonant Inventory}

The consonant phonemes may be organized by place and manner of articulation as follows.

\begin{center}
\begin{tabular}{lcccccc}
 & Bilabial & Labiodental & Alveolar & Postalveolar & Velar & Glottal \\ \hline
Plosive (voiceless) & p &  & t &  & k &  \\
Plosive (voiced) & b &  & d &  & g &  \\
Fricative (voiceless) &  & f & s & ʃ &  & h \\
Fricative (voiced) &  & v & z &  &  &  \\
Affricate &  &  &  & tʃ &  &  \\
Nasal & m &  & n &  & ŋ &  \\
Liquid &  &  & l, r &  &  &  \\
Glide & w &  &  &  & j &  \\
\end{tabular}
\end{center}

\subsection{Observations on Consonant Structure}

Consonant clusters are limited. Word-initial clusters are generally restricted to those inherited from English borrowings (e.g., \textit{school}, \textit{small}), though epenthesis or reduction frequently occurs in casual speech. Word-final obstruents are typically unreleased and may undergo devoicing.

The rhotic phoneme exhibits variable realization, ranging from alveolar tap to approximant depending on speaker background.

The velar nasal /ŋ/ occurs syllable-finally and in medial position but rarely word-initially.

\subsection{Vowel Inventory}

The vowel system is comparatively reduced relative to Standard English. A five- to seven-vowel analysis is common. A conservative six-vowel system is given below.

\begin{center}
\begin{tabular}{lccc}
 & Front & Central & Back \\ \hline
High & i &  & u \\
Mid & e &  & o \\
Low &  & a &  \\
\end{tabular}
\end{center}

Diphthongs such as /ai/, /au/, and /oi/ occur primarily in English-derived lexemes.

\subsection{Prosodic Structure}

Nigerian Pidgin is not tonally contrastive in the phonemic sense typical of many substrate languages. However, intonation patterns reflect substrate influence and may encode pragmatic distinctions.

Stress placement is typically inherited from English source forms but tends toward simplification in polysyllabic borrowings.

\subsection{Phonological Implications for Orthography}

The phonological inventory imposes a set of structural constraints that are directly relevant to orthographic design. First, the vowel system is comparatively small and exhibits a high degree of stability across regional varieties. This relative simplicity renders consistent vowel marking both feasible and efficient, as the number of contrasts to be represented does not generate excessive diacritic load. An orthography may therefore encode vocalic distinctions explicitly without introducing undue graphic complexity.

Second, while most consonantal distinctions are inherited from English, the language does not productively generate complex consonant clusters beyond those preserved in borrowings. The limited productivity of clustering reduces the need for elaborate graphemic mechanisms to encode multi-consonantal sequences. Orthographic representation can thus remain structurally transparent, with minimal accommodation for internally generated cluster variation.

Third, the absence of phonemic tone significantly reduces orthographic burden when compared with tonal languages such as Yoruba. Because tonal contrast does not function to distinguish lexical meaning in Nigerian Pidgin, the script need not incorporate systematic tonal marking. This constraint simplifies graphic design and enhances legibility, particularly in contexts where diacritic density might otherwise hinder adoption.

Taken together, these phonological properties support the viability of a streamlined and internally coherent orthographic system. The structural profile of the language aligns with a script architecture that prioritizes phonetic transparency, limited diacritic complexity, and restrained graphemic elaboration.

Any orthographic system must therefore preserve:

\[
\text{Contrast}(p,b),\;
\text{Contrast}(t,d),\;
\text{Contrast}(k,g),\;
\text{Contrast}(s,ʃ),\;
\text{Contrast}(f,v)
\]

while maintaining vowel distinctiveness sufficient to avoid lexical ambiguity.

\subsection{Conclusion}

This phoneme inventory provides the foundational layer upon which the proposed Arabic-script orthography is constructed. Orthographic evaluation must proceed downstream from phonological specification. Without this inventory, grapheme selection cannot be systematically justified.

\section{Appendix B: Grapheme--Phoneme Correspondence Table}

\subsection{Principles of Mapping}

The proposed orthography restricts itself to the Classical Arabic consonant inventory and uses a small number of explicit digraph conventions to represent Nigerian Pidgin phonemes not encoded as single Arabic letters. The aim is not perfect reversibility in the absence of vowel marks, but a stable pedagogical correspondence that can be made fully deterministic under full vowel marking.

Formally, the system defines a mapping
\[
\text{(phoneme or phoneme cluster)} \longrightarrow \text{(Arabic grapheme sequence)}
\]
with the convention that certain sequences of letters are treated as single units for reading purposes.

\subsection{Consonant Correspondence}

\begin{center}
\begin{tabular}{lll}
\textbf{Phoneme (IPA)} & \textbf{Arabic Grapheme(s)} & \textbf{Example (Pidgin)} \\ \hline
/p/ & \AR{ب} & \AR{بِيجِنْ} (pidgin) \\
/b/ & \AR{ب} & \AR{بُوِي} (boy) \\
/t/ & \AR{ت} & \AR{تِيكْ} (take) \\
/d/ & \AR{د} & \AR{دِمْ} (dem) \\
/k/ & \AR{ك} & \AR{كَمْ} (come) \\
/g/ & \AR{ك}\,(context) & \AR{كُو} (go) \\
/f/ & \AR{ف} & \AR{فِيتْ} (fit) \\
/v/ & \AR{ف} & \AR{فِرِي} (very) \\
/s/ & \AR{س} & \AR{سِي} (see) \\
/\textipa{\textesh}/ & \AR{ش} & \AR{شُورْ} (sure) \\
/z/ & \AR{ز} & \AR{زُو} (zoo) \\
/t\textipa{\textesh}/ & \AR{ت\,+\,ش} & \AR{تْشُوبْ} (chop) \\
/m/ & \AR{م} & \AR{مِي} (me) \\
/n/ & \AR{ن} & \AR{نُو} (no) \\
/\textipa{N}/ & \AR{ن}\,(context) & \AR{سِنْ} (sing) \\
/l/ & \AR{ل} & \AR{لايْفْ} (life) \\
/r/ & \AR{ر} & \AR{رَنْ} (run) \\
/w/ & \AR{و} & \AR{وِي} (we) \\
/j/ & \AR{ي} & \AR{يَانْ} (yarn) \\
/h/ & \AR{ه} & \AR{هِيَرْ} (hear) \\
\end{tabular}
\end{center}

The system intentionally merges certain contrasts that are not represented in Classical Arabic as single letters. In particular, /p/ is written with \AR{ب} and /v/ is written with \AR{ف}. The affricate /t\textipa{\textesh}/ is written as the explicit digraph \AR{ت+ش}, which is interpreted as a single consonant in reading.

\subsection{The \texorpdfstring{\AR{ت+ش}}{t+sh} Digraph Convention}

Let the digraph operator be defined by
\[
\mathrm{CH} := \AR{ت}\,\AR{ش}.
\]
Then the reading rule is:
\[
\mathrm{CH} \mapsto /t\textipa{\textesh}/
\]
regardless of whether the two letters could otherwise be parsed separately. When desired for pedagogical clarity, a suk\=un may be placed on \AR{ت} to signal clustering:
\[
\AR{تْش} \quad \text{rather than} \quad \AR{تش}.
\]
No additional letters beyond the Classical inventory are required.

\subsection{Vowel Correspondence}

The orthography permits full vowel marking in formal or instructional contexts and partial vowel omission in fluent contexts. A pedagogical one-to-one mapping is assumed when vowels are marked.

\begin{center}
\begin{tabular}{lll}
\textbf{Phoneme (IPA)} & \textbf{Arabic Representation} & \textbf{Example} \\ \hline
/i/ & \AR{ِ} (kasra) or \AR{ي} & \AR{سِي} (see) \\
/e/ & \AR{ِ} or \AR{ِ} \AR{يْ} & \AR{بِلِيفْ} (believe) \\
/a/ & \AR{َ} (fatha) or \AR{ا} & \AR{مَاكْ} (make) \\
/o/ & \AR{ُ} (damma) or \AR{و} & \AR{كُو} (go) \\
/u/ & \AR{ُ} or \AR{و} & \AR{فُودْ} (food) \\
/ai/ & \AR{اي} & \AR{لايْفْ} (life) \\
/au/ & \AR{او} & \AR{ناو} (now) \\
/oi/ & \AR{اوي} & \AR{بُوِي} (boy) \\
\end{tabular}
\end{center}


\subsection{Orthographic Determinism Under Full Marking}

Let $\Sigma_P$ denote the intended Pidgin phoneme inventory and let $\Sigma_A^\ast$ denote finite Arabic grapheme strings. Define the encoding map
\[
G : \Sigma_P \rightarrow \Sigma_A^\ast
\]
with $G(t\textipa{\textesh})=\text{ت}\text{ش}$ as a primitive assignment. Under full vowel marking, the reading map becomes operationally deterministic for the intended learner audience. When vowels are omitted, ambiguity is treated as a controlled property of the script rather than a defect, paralleling ordinary Arabic orthographic practice.

\subsection{Cluster Representation}

Consonant clusters are written sequentially. Optional suk\=un marking may be used to prevent unintended vowel epenthesis in novice reading. No epenthetic vowels are introduced orthographically unless they are phonologically realized in the target utterance.

\subsection{Conclusion}

This correspondence table defines a Classical-Arabic-only Ajami-style orthography for Nigerian Pidgin with a minimal digraph extension for /t\textipa{\textesh}/. The result preserves typographic accessibility while keeping the mapping rules explicit and formally describable.

\section{Appendix C: Syllable Structure and Phonotactic Constraints}

\subsection{Syllable Template}

The syllable structure of Nigerian Pidgin can be characterized by the general template

[
\sigma = (C)(C)V(C),
]

in which ( C ) represents a consonantal segment and ( V ) denotes the vocalic nucleus. Parentheses mark optional positions, indicating that both single and complex onsets are permitted, while codas remain structurally constrained.

Empirically, the most frequent syllable configurations are ( CV ), ( CVC ), ( V ), and ( CCV ). The dominance of open syllables and simple codas reflects a phonotactic profile that favors clarity of articulation and rhythmic regularity. Complex codas are comparatively rare and are typically traceable to English lexical inheritance rather than to internally generated phonological processes. In many cases, such clusters undergo simplification in natural speech, further reinforcing the preference for less structurally dense syllable margins.

\subsection{Onset Structure}

Single-consonant onsets constitute the unmarked case. Forms such as *go* instantiate the ( CV ) pattern, while items like *fit* exemplify the ( CVC ) structure, demonstrating the compatibility of simple codas with otherwise minimal syllable architecture.

Clustered onsets, by contrast, appear predominantly in borrowings from English. Lexical items such as *small* and *school* exhibit initial consonant clusters corresponding to ( CCVC ) patterns. However, even in these cases, phonological reduction is common in rapid or casual speech. The cluster may be compressed or reanalyzed prosodically, producing realizations such as [smɔl] or a form with reduced initial prominence. This variability indicates that while inherited clusters are tolerated, they are not strongly reinforced by native phonological processes.

Importantly, there is no evidence of productive cluster formation independent of English-derived vocabulary. Consonant clustering in Nigerian Pidgin is therefore best understood as lexically inherited rather than as a generative phonotactic strategy internal to the language.

\subsection{Coda Structure}

Word-final codas are typically single consonants:

\[
\text{fit} \rightarrow CVC
\]
\[
\text{run} \rightarrow CVC
\]

Complex codas such as /-nd/ or /-st/ are possible but are variably simplified:

\[
\text{last} \rightarrow [las] \text{ in casual speech}
\]

Thus coda clusters are lexically inherited rather than generatively productive.

\subsection{Vowel Reduction and Stability}

Unlike Standard English, vowel reduction in unstressed syllables is minimal. The system tends toward vowel preservation rather than centralized schwa collapse.

This relative vowel stability supports explicit vowel marking in orthography without excessive burden.

\subsection{Phonotactic Constraints}

The following constraints are broadly observed:

\[
\text{No word-initial } /ŋ/
\]

\[
\text{No triple consonant clusters}
\]

\[
\text{Limited coda complexity}
\]

\[
\text{Preference for open syllables in non-borrowed lexicon}
\]

\subsection{Implications for Arabic-Script Adaptation}

The phonotactic profile supports direct consonant sequencing without epenthetic insertion. Since clusters are limited and predictable, the orthography does not require additional structural markers beyond linear concatenation.

Let a Pidgin word be represented as:

\[
W = C_1 C_2 V C_3
\]

Orthographic representation preserves cluster integrity:

\[
C_1 C_2 V C_3 \rightarrow G(C_1)G(C_2)G(V)G(C_3)
\]

No morphophonemic alternations require special graphemic treatment.

\subsection{Formal Constraint Statement}

Define allowable syllable set \( \mathcal{S} \) such that:

\[
\mathcal{S} = \{ V, CV, CVC, CCV, CCVC \}
\]

with CCV and CCVC restricted to lexically inherited forms.

Orthographic validity requires:

\[
\forall \sigma \in \mathcal{S}, \; G(\sigma) \text{ is representable without additional diacritics beyond vowel marking.}
\]

\subsection{Conclusion}

The phonotactic system of Nigerian Pidgin is structurally simple relative to English and substrate languages. This simplicity facilitates orthographic determinism and minimizes ambiguity in Arabic-script adaptation.

\section{Appendix D: Morphosyntactic Structure and Functional Tagging}

\subsection{Overview}

This appendix provides a formal description of core morphosyntactic operators attested in the corpus. Nigerian Pidgin is largely analytic, with grammatical relations expressed through particles rather than inflectional morphology.

We adopt a functional notation in which lexical roots are represented as \( R \), grammatical particles as \( G \), and argument positions as \( A_i \).

\subsection{Tense–Aspect–Mood System}

The tense–aspect system is particle-based. Let \( V \) denote a verb root.

\subsubsection*{Progressive Marker \textit{dey}}

The progressive construction may be formalized as:

\[
S \rightarrow NP + \text{dey} + V
\]

Example:

\[
\text{I dey waka}
\]

\[
NP_1 + G_{PROG} + V
\]

Here, \textit{dey} encodes imperfective progressive aspect.

\subsubsection*{Perfective Marker \textit{don}}

The perfective construction may be formalized:

\[
S \rightarrow NP + \text{don} + V
\]

Example:

\[
\text{Dem don come}
\]

\[
NP_1 + G_{PERF} + V
\]

\textit{don} marks completed action without inflectional agreement.

\subsubsection*{Modal Marker \textit{fit}}

Ability or possibility is expressed via:

\[
S \rightarrow NP + \text{fit} + V
\]

Example:

\[
\text{You fit do am}
\]

This corresponds to modal \( CAN \).

\subsection{Negation}

Negation is pre-verbal and invariant:

\[
S \rightarrow NP + \text{no} + V
\]

Example:

\[
\text{I no sabi}
\]

There is no auxiliary inversion or do-support.

\subsection{Copular Constructions}

Two copular strategies exist.

\subsubsection*{Equative Copula \textit{na}}

\[
S \rightarrow \text{na} + NP
\]

Example:

\[
\text{Na true}
\]

\textit{na} marks equative identification.

\subsubsection*{Zero Copula}

Predicate adjectives often appear without overt copula:

\[
\text{She no well}
\]

This corresponds structurally to:

\[
NP + NEG + ADJ
\]

\subsection{Pronominal System}

The core pronoun inventory may be represented:

\begin{center}
\begin{tabular}{ll}
\textbf{Function} & \textbf{Form} \\ \hline
1SG & I \\
2SG & you \\
3SG & e / she / he \\
1PL & we \\
2PL & una \\
3PL & dem \\
\end{tabular}
\end{center}

Pronouns do not inflect for case.

\subsection{Object Marker \textit{am}}

Third-person objects frequently appear as invariant particle:

\[
S \rightarrow NP_1 + V + am
\]

Example:

\[
\text{I do am}
\]

This functions as a cliticized object pronoun.

\subsection{Serial Verb Construction}

Serial verb constructions occur without conjunction:

\[
NP + V_1 + V_2
\]

Example:

\[
\text{Make we go see}
\]

This reflects substrate influence from West African languages.

\subsection{Clause Typology}

Interrogatives are formed via intonation or question words without inversion:

\[
\text{You dey come?}
\]

No auxiliary movement occurs.

\subsection{Formal Summary}

Let:

\[
\mathcal{G} = \{\text{dey}, \text{don}, \text{fit}, \text{no}, \text{na}\}
\]

The grammar is characterized by:

\[
\text{No inflectional morphology}
\]

\[
\text{Pre-verbal TAM particles}
\]

\[
\text{Invariant negation}
\]

\[
\text{Minimal agreement marking}
\]

Thus Nigerian Pidgin may be modeled as an analytic grammar:

\[
S = NP + G^* + V + (NP)
\]

where \( G^* \) denotes zero or more grammatical particles.

\subsection{Implications for Orthography}

Because morphosyntactic distinctions are particle-based and not inflectional, orthographic clarity depends on stable representation of these particles:

\[
\text{dey}, \text{don}, \text{fit}, \text{no}, \text{na}
\]

Any ambiguity in their graphemic rendering would introduce grammatical ambiguity.

\subsection{Conclusion}

The morphosyntactic structure of Nigerian Pidgin is typologically analytic, particle-driven, and structurally transparent. Orthographic stability therefore depends less on morphological marking and more on consistent representation of invariant grammatical operators.

\section{Appendix E: Comparative Analysis with Hausa Ajami and West African Arabic-Script Traditions}

\subsection{Scope and Purpose}

This appendix situates the proposed Nigerian Pidgin Arabic-script system within the broader West African Ajami tradition. The objective is structural comparison rather than historical narrative.

Ajami refers to the adaptation of the Arabic script to represent non-Arabic languages. Hausa Ajami provides the most developed and standardized example in West Africa and therefore serves as the principal comparative baseline.

![Image](https://www.csmc.uni-hamburg.de/14438068/ajami-lab-02-48385be4f7cf110a0021950e1430417e5d3c2e0c.jpg)

![Image](https://i.etsystatic.com/32646785/r/il/21c14b/7223405767/il_570xN.7223405767_4pw0.jpg)

![Image](https://r12a.github.io/scripts/arab/ha/fig_naskh_style.png)

![Image](https://www.omniglot.com/images/writing/hausa_aj.gif)

\subsection{Hausa Ajami: Structural Features}

Hausa Ajami represents a systematic adaptation of the Arabic script to the phonological structure of Hausa. Its development illustrates how an inherited consonantal script can be expanded to accommodate a distinct sound system through principled modification rather than wholesale replacement.

One central strategy involves the extension of the grapheme inventory through modified letters. Additional characters such as پ for /p/ and گ for /g/ exemplify the introduction of diacritic or structural variation to encode consonants absent in Classical Arabic. This approach preserves visual continuity with the Arabic script while increasing phonemic coverage.

A second strategy concerns vowel representation. Hausa Ajami employs systematic vowel marking, either through diacritics or through the use of long-vowel letters functioning as matres lectionis. Given the importance of vowel quality in Hausa lexical contrast, such marking practices enhance phonological transparency and reduce ambiguity in written form.

Tone, which is phonemic in Hausa, is not ordinarily marked explicitly in Ajami writing. Instead, tonal interpretation is typically inferred from lexical familiarity and syntactic context. This indirect accommodation of tone reflects a pragmatic balance between phonological precision and orthographic economy.

Hausa phonology includes implosives, ejectives, and tonal contrasts that are not present in Nigerian Pidgin. As a result, Hausa Ajami necessitates a greater degree of graphemic innovation to encode these distinctions. The comparison underscores that orthographic adaptation is conditioned by phonological complexity: where the sound system demands fine-grained contrast, the script must evolve correspondingly.

Formally, Hausa mapping may be expressed:

\[
G_H : \Sigma_{Hausa} \rightarrow \Sigma_{Ajami}
\]

with partial context sensitivity due to tone omission.

\subsection{Structural Comparison with Nigerian Pidgin}

Let:

\[
\Sigma_P = \text{Pidgin phoneme set}
\]
\[
\Sigma_H = \text{Hausa phoneme set}
\]

Then:

\[
|\Sigma_P| < |\Sigma_H|
\]

Pidgin lacks:

\[
\text{Tone contrast}
\]
\[
\text{Implosive consonants}
\]
\[
\text{Ejective consonants}
\]

Therefore, the orthographic burden required for Pidgin Ajami is strictly lower.

\subsection{Graphemic Extension Parallels}

Both Hausa Ajami and the proposed Pidgin system rely on extended Arabic letters:

\[
پ, گ, چ, ڤ
\]

These are not foreign innovations but established orthographic tools within Islamic manuscript traditions.

Thus the Pidgin proposal does not introduce structurally novel graphemes relative to regional precedent.

\subsection{Vowel Representation}

Hausa Ajami often under-specifies short vowels in non-pedagogical texts, relying on reader inference.

The proposed Pidgin system recommends optional full vowel marking in educational contexts and partial reduction in fluent writing.

This mirrors established Ajami practice:

\[
\text{Full marking} \rightarrow \text{pedagogical precision}
\]
\[
\text{Reduced marking} \rightarrow \text{scribal economy}
\]

\subsection{Morphological Transparency}

Hausa is morphologically richer than Pidgin. Hausa Ajami must encode:

\[
\text{Gender}
\]
\[
\text{Number agreement}
\]
\[
\text{Derivational morphology}
\]

Nigerian Pidgin, being analytic, imposes fewer morphological marking requirements.

Thus orthographic transparency may be achieved with lower complexity.

\subsection{Script Directionality and Cultural Integration}

Ajami scripts are written right-to-left. Adoption of Arabic script for Pidgin would therefore:

\[
\text{Align with established regional literacy traditions}
\]

However, unlike Hausa, Nigerian Pidgin has not historically been transmitted through Islamic scholarly networks. Therefore sociolinguistic adoption dynamics would differ.

\subsection{Formal Comparative Statement}

Let orthographic complexity be approximated as:

\[
C = |\Sigma_G| + D + T
\]

where:

\[
|\Sigma_G| = \text{number of graphemes}
\]
\[
D = \text{diacritic burden}
\]
\[
T = \text{tone marking complexity}
\]

For Hausa Ajami:

\[
T > 0
\]

For Nigerian Pidgin Ajami:

\[
T = 0
\]

Therefore:

\[
C_{Pidgin} < C_{Hausa}
\]

under equivalent vowel-marking regimes.

\subsection{Conclusion}

The proposed Pidgin Arabic orthography does not represent an unprecedented script innovation. It operates within established Ajami structural patterns while benefiting from reduced phonological complexity relative to Hausa.

The principal challenge is not graphemic feasibility but sociolinguistic adoption.

\section{Appendix F: Orthographic Economy and Information Density}

\subsection{Objective}

This appendix evaluates the proposed Arabic-script orthography for Nigerian Pidgin in terms of orthographic economy and information density. The aim is to quantify structural efficiency rather than to argue aesthetically.

Orthographic economy is defined as the ratio between phonological information encoded and graphemic length required.

\subsection{Definitions}

Let:

\[
W_L = \text{Latin-script representation length}
\]

\[
W_A = \text{Arabic-script representation length}
\]

Let \( |W| \) denote grapheme count excluding whitespace.

Define the compression ratio:

\[
R = \frac{|W_A|}{|W_L|}
\]

\subsection{Empirical Sample}

Consider representative examples from the corpus.

\[
\text{I dey come}
\]

Latin length:

\[
|W_L| = 9
\]

Arabic:

\[
\AR{آيْ دي كَمْ}
\]

\[
|W_A| = 8
\]

Thus:

\[
R = \frac{8}{9} \approx 0.89
\]

Second example:

\[
\text{You no get sense}
\]

Latin:

\[
|W_L| = 14
\]

Arabic:

\[
\AR{يُو نَوْ جِتْ سِنْسْ}
\]

\[
|W_A| = 13
\]

\[
R \approx 0.93
\]

Across sampled entries, preliminary averaging yields:

\[
R \in [0.85, 1.05]
\]

indicating near parity in graphemic density.

\subsection{Vowel Marking and Density}

If full diacritic marking is retained, character count increases marginally.

If short vowels are omitted in fluent writing, the ratio decreases:

\[
R_{\text{reduced}} < R_{\text{full}}
\]

Thus vowel omission provides optional compression at the cost of increased reliance on contextual disambiguation.

\subsection{Information-Theoretic Framing}

Let:

\[
H(P) = -\sum p(x)\log p(x)
\]

represent entropy of phoneme distribution.

Orthographic efficiency may be approximated as:

\[
E = \frac{H(P)}{\text{Average grapheme length}}
\]

Since Pidgin has relatively low phonemic entropy due to reduced vowel inventory and limited cluster productivity, the system supports compact encoding.

\subsection{Redundancy Analysis}

Latin orthography for Pidgin inherits English irregularities, including silent letters and digraph ambiguity:

\[
\text{“ough” phenomena absent but analogous ambiguity exists}
\]

Arabic-script representation is phonemic and thus reduces redundancy.

Redundancy \( \rho \) may be approximated as:

\[
\rho = 1 - \frac{\text{phoneme count}}{\text{grapheme count}}
\]

The proposed system aims to minimize \( \rho \).

\subsection{Comparative Economy}

Let:

\[
C_L = \text{Latin complexity}
\]
\[
C_A = \text{Arabic complexity}
\]

If Latin representation is quasi-etymological while Arabic representation is phonemic, then:

\[
C_A \leq C_L
\]

under consistent vowel marking conventions.

\subsection{Cognitive Load Considerations}

Right-to-left directionality introduces initial learning cost. However, grapheme inventory size remains modest:

\[
|\Sigma_A| \approx 29
\]

which lies within standard literacy acquisition norms.

\subsection{Conclusion}

Quantitative evaluation indicates that the proposed orthographic system maintains near parity in graphemic length when compared with its Latin-based counterpart. The average number of written symbols required to encode a lexical item does not increase in any systematic way, suggesting that the transition in script does not entail expansion at the level of surface representation.

At the same time, the system reduces phoneme-to-grapheme ambiguity by tightening the correspondence between sound and symbol. Because the foundational layer encodes phonological structure directly, it avoids many of the irregularities characteristic of conventional English orthography. This increased transparency enhances reversibility and predictability, particularly at the stage of phonetic reconstruction.

The framework also permits optional compression through the omission of diacritics in contexts where ambiguity is unlikely. Such flexibility allows writers to modulate precision and density according to communicative needs, preserving structural integrity while accommodating practical usage.

From an information-density perspective, the system is therefore structurally efficient. It does not introduce measurable expansion relative to Latin representation, and in certain contexts may even reduce interpretive entropy by aligning graphemic structure more closely with phonological form.

\section{Appendix G: Ambiguity and Error Propagation Analysis}

\subsection{Objective}

This appendix evaluates potential sources of ambiguity and error propagation within the proposed Arabic-script orthography for Nigerian Pidgin. The goal is to identify structural weaknesses and determine whether they arise from the phonology of the language or from script adaptation.

\subsection{Formal Model of Orthographic Mapping}

Let:

\[
G : \Sigma_P^* \rightarrow \Sigma_A^*
\]

where \( \Sigma_P \) is the phoneme inventory and \( \Sigma_A \) the grapheme inventory.

Ambiguity arises if:

\[
\exists x,y \in \Sigma_P^*, \; x \neq y \;\text{such that}\; G(x) = G(y)
\]

We analyze potential cases.

\subsection{Vowel Omission Ambiguity}

If short vowels are omitted, sequences such as:

\[
\text{fit} \quad \text{and} \quad \text{fat}
\]

could reduce to identical consonantal skeletons.

Example:

\[
فْتْ
\]

Thus vowel omission introduces recoverable but real ambiguity.

Mitigation strategy: mandatory vowel marking in minimal-pair environments.

\subsection{Homophony Inherited from Pidgin}

Certain lexical items are homophonous independent of script.

For example:

\[
\text{no} \quad (\text{negation}) 
\]
\[
\text{know} \quad (\text{lexical verb})
\]

In Pidgin these collapse phonologically:

\[
\AR{نو}
\]

The ambiguity is phonemic rather than orthographic. The script therefore does not introduce instability beyond that already present in spoken language.

\subsection{Consonant Cluster Compression}

Clusters such as:

\[
\AR{سْمُولْ}
\]

may appear visually dense but do not collapse phonemic distinction. No graphemic ambiguity arises provided cluster letters remain distinct.

\subsection{Loanword Adaptation}

English-derived terms may exhibit variable vowel realization. For example:

\[
\text{school} \rightarrow \AR{سكول}
\]

If speakers variably pronounce epenthetic vowels, orthographic normalization may diverge from speech variation. This represents a standardization tension rather than a mapping ambiguity.

\subsection{Diacritic Loss in Informal Writing}

If diacritics are omitted entirely:

\[
\AR{كتب}
\]

ambiguity increases substantially, since short vowels are no longer recoverable from graphemic form alone.

Define diacritic omission function:

\[
D : \Sigma_A \rightarrow \Sigma_A'
\]

where \( \Sigma_A' \subset \Sigma_A \) excludes short vowel marks.

Then ambiguity probability increases as:

\[
P(\text{collision}) \propto \frac{|\Sigma_P|}{|\Sigma_A'|}
\]

Full vowel marking therefore reduces collision probability.

\subsection{Error Propagation in Iterative Copying}

Consider iterative transcription:

\[
W_0 \rightarrow W_1 \rightarrow \dots \rightarrow W_n
\]

If each stage introduces independent error probability \( \epsilon \), then expected distortion after \( n \) iterations is:

\[
E[D_n] = 1 - (1-\epsilon)^n
\]

Phonemic orthography reduces \( \epsilon \) relative to quasi-etymological systems, since fewer silent letters or irregular mappings exist.

\subsection{Structural Stability Conditions}

The orthography remains injective under the following constraints:

\[
\text{Full vowel marking}
\]
\[
\text{Retention of extended graphemes}
\]
\[
\text{No grapheme conflation}
\]

If these conditions hold, then:

\[
G \text{ is injective}
\]

If diacritics are omitted:

\[
G \text{ becomes partially many-to-one}
\]

\subsection{Comparative Ambiguity}

Latin-script Pidgin already contains ambiguity due to English spelling conventions. For example:

\[
\text{read} \; (\text{present}) \neq \text{read} \; (\text{past})
\]

Such orthographic instability is inherited from English rather than intrinsic to Pidgin.

The Arabic-script proposal, being phonemic, may reduce certain inherited ambiguities.

\subsection{Conclusion}

Ambiguity in the proposed system arises primarily from optional vowel omission rather than from consonant representation. When vowels are fully marked, the mapping is structurally robust and minimally ambiguous.

Error propagation rates are comparable to, and potentially lower than, Latin-script representation due to phonemic transparency.

\section{Appendix H: Reversibility Theorem and Transliteration Functions}

\subsection{Definitions}

Let:

\[
\Sigma_P = \text{Pidgin phoneme inventory}
\]

\[
\Sigma_A = \text{Arabic grapheme inventory}
\]

\[
\Sigma_L = \text{Latin grapheme inventory}
\]

Define transliteration functions:

\[
G : \Sigma_P^* \rightarrow \Sigma_A^*
\]

\[
L : \Sigma_P^* \rightarrow \Sigma_L^*
\]

and inverse mappings:

\[
G^{-1} : \Sigma_A^* \rightarrow \Sigma_P^*
\]

\[
L^{-1} : \Sigma_L^* \rightarrow \Sigma_P^*
\]

\subsection{Orthographic Determinism}

The proposed system assigns a unique grapheme to each phoneme:

\[
\forall p \in \Sigma_P,\; G(p) = a \in \Sigma_A
\]

such that:

\[
p_1 \neq p_2 \Rightarrow G(p_1) \neq G(p_2)
\]

Thus:

\[
G \text{ is injective at the phoneme level.}
\]

\subsection{Weak Reversibility Theorem}

\textbf{Theorem.}  
If all vowel diacritics are preserved, then:

\[
G^{-1} \circ G = \mathrm{Id}_{\Sigma_P^*}
\]

That is, transliteration to Arabic script and back to phonemic representation recovers the original sequence.

\textbf{Proof.}

Since \( G \) is injective over \( \Sigma_P \), and since concatenation preserves symbol boundaries, we have:

\[
G(p_1 p_2 \dots p_n) = G(p_1)G(p_2)\dots G(p_n)
\]

Because each \( G(p_i) \) corresponds uniquely to \( p_i \), the inverse mapping recovers:

\[
G^{-1}(G(p_i)) = p_i
\]

Thus:

\[
G^{-1}(G(p_1 p_2 \dots p_n)) = p_1 p_2 \dots p_n
\]

provided no grapheme ambiguity is introduced by diacritic omission.

\hfill \( \square \)

\subsection{Partial Reversibility Under Vowel Omission}

Let \( D \) be a diacritic removal function:

\[
D : \Sigma_A^* \rightarrow \Sigma_A'^*
\]

where \( \Sigma_A' \subset \Sigma_A \) excludes short vowel markers.

Then composite mapping:

\[
G' = D \circ G
\]

may no longer be injective.

\textbf{Corollary.}  
If two phonemic sequences differ only in vowel quality and short vowels are unmarked, then:

\[
G'(x) = G'(y)
\]

even if \( x \neq y \).

Thus reversibility becomes context-dependent rather than absolute.

\subsection{Latin Transliteration Comparison}

Latin-script Pidgin is not strictly phonemic. Let:

\[
L(p) = \ell
\]

English orthographic conventions introduce non-phonemic irregularities. Therefore:

\[
L^{-1} \text{ is not strictly deterministic without phonological knowledge.}
\]

Hence, under full vowel marking:

\[
\text{Arabic-script representation may exhibit stronger phonemic reversibility than Latin representation.}
\]

\subsection{Algorithmic Implementation}

Define transliteration procedure:

\[
T_A: \text{Input Latin Pidgin} \rightarrow \text{Phonemic parse} \rightarrow G(\cdot)
\]

\[
T_L: \text{Input Arabic Pidgin} \rightarrow G^{-1}(\cdot) \rightarrow \text{Latin rendering}
\]

Since the mapping operates at phoneme level rather than orthographic level, intermediate phonemic parsing ensures structural consistency.

\subsection{Complexity}

Let word length be \( n \).

Then transliteration complexity:

\[
O(n)
\]

as each symbol is mapped independently.

No recursive or context-sensitive rewriting is required under full marking.

\subsection{Conclusion}

The proposed orthography satisfies:

\[
G^{-1} \circ G = \mathrm{Id}
\]

under full vowel marking.

Under vowel reduction, reversibility is weakened but remains recoverable via lexical context.

Thus the system is formally well-defined, computationally tractable, and structurally reversible within specified constraints.

\section{Appendix I: Minimal Pair Dataset and Contrast Preservation}

\subsection{Objective}

This appendix evaluates whether the proposed Arabic-script orthography preserves phonemic contrasts attested in Nigerian Pidgin. Orthographic adequacy requires that minimal phonemic distinctions correspond to distinct graphemic forms.

Formally, if:

\[
x, y \in \Sigma_P^*, \quad x \neq y
\]

and

\[
\text{Contrast}(x,y) \text{ is phonemic}
\]

then it must hold that:

\[
G(x) \neq G(y)
\]

\subsection{Consonant Contrast Tests}

\subsubsection*{Plosive Voicing Contrast}

\[
\text{pat} \neq \text{bat}
\]

\[
\AR{پَتْ} \neq \AR{بَتْ}
\]

Distinct graphemes preserve contrast.

\subsubsection*{Velar Contrast}

\[
\text{go} \neq \text{ko}
\]

\[
\AR{گُو} \neq \AR{كُو}
\]

Contrast is preserved.

\subsubsection*{Fricative Contrast}

\[
\text{fan} \neq \text{van}
\]

\[
\AR{فَن} \neq \AR{ڤَن}
\]

Distinct graphemes maintain phonemic distinction.

\subsubsection*{Sibilant Contrast}

\[
\text{see} \neq \text{she}
\]

\[
\AR{سِي} \neq \AR{شِي}
\]

Contrast is preserved.

\subsection{Vowel Contrast Tests}

\subsubsection*{/i/ vs /a/}

\[
\text{fit} \neq \text{fat}
\]

\[
\AR{فِيتْ} \neq \AR{فَتْ}
\]

Distinct vowel marking is required for full preservation.

\subsubsection*{/o/ vs /u/}

\[
\text{cot} \neq \text{cut}
\]

\[
\AR{كُتْ} \neq \AR{كَتْ}
\]

Vowel marking disambiguates.

\subsection{Modal vs Lexical Distinction}

\[
\text{fit} \; (\text{modal}) \neq \text{fit} \; (\text{adjective})
\]

These forms are homophonous in Pidgin and therefore not orthographically distinguishable in any phonemic system.

Thus:

\[
\text{Ambiguity inherited from phonology}
\]

rather than from script design.

\subsection{Particle Contrast Preservation}

\[
\text{no} \neq \text{don}
\]

\[
\AR{نو} \neq \AR{دُونْ}
\]

Distinct grammatical operators remain visually distinct.

\subsection{Cluster Preservation}

\[
\text{small} \neq \text{mall}
\]

\[
\AR{سْمُولْ} \neq \AR{مُولْ}
\]

The initial cluster is preserved without conflation.


\subsection{Formal Statement}

Let:

\[
\mathcal{M} = \{(x,y) \mid x,y \text{ minimal pairs}\}
\]

The orthography satisfies:

\[
\forall (x,y) \in \mathcal{M}, \quad G(x) \neq G(y)
\]

under full vowel marking.

If vowel marking is omitted:

\[
\exists (x,y) \in \mathcal{M} \text{ such that } G'(x) = G'(y)
\]

only when the contrast depends solely on short vowels.

\subsection{Contrast Preservation Theorem}

\textbf{Theorem.}  
Under full vowel marking, the proposed orthography preserves all consonantal and vocalic phonemic contrasts present in the defined phoneme inventory.

\textbf{Proof Sketch.}

Since mapping \( G \) is injective at phoneme level and since minimal pairs differ in at least one phoneme, their grapheme sequences must differ in at least one position.

\[
x \neq y \Rightarrow \exists i \; (p_i \neq q_i)
\]

\[
G(p_i) \neq G(q_i)
\]

Thus:

\[
G(x) \neq G(y)
\]

\hfill \( \square \)

\subsection{Conclusion}

The orthography preserves phonemic contrast under full specification. Ambiguities arise only under deliberate vowel omission and are therefore optional rather than structural failures.

\section{Appendix J: Computational Transliteration Specification Under Classical Arabic Constraints}

\subsection{Alphabet Restriction}

The Arabic grapheme inventory is restricted to:

\[
\Sigma_A^{\text{classical}}
=
\{
\AR{ا}, \AR{ب}, \AR{ت}, \AR{ث}, \AR{ج}, \AR{ح}, \AR{خ}, \AR{د}, \AR{ذ},
\AR{ر}, \AR{ز}, \AR{س}, \AR{ش}, \AR{ص}, \AR{ض}, \AR{ط}, \AR{ظ},
\AR{ع}, \AR{غ}, \AR{ف}, \AR{ق}, \AR{ك}, \AR{ل}, \AR{م}, \AR{ن},
\AR{ه}, \AR{و}, \AR{ي}
\}
\]

No Persian-derived letters (\AR{پ}, \AR{گ}, \AR{چ}, \AR{ڤ}) are permitted.

\subsection{Phoneme Set}

Let:

\[
\Sigma_P = \text{Pidgin phoneme inventory defined in Appendix A}
\]

including phonemes absent in Classical Arabic:

\[
\{ p, g, v, tʃ \}
\]

\subsection{Digraph Encoding Strategy}

Since injective single-grapheme mapping is impossible for these phonemes, we define a sequence-level mapping:

\[
G : \Sigma_P \rightarrow (\Sigma_A^{classical})^+
\]

where the codomain permits multi-grapheme outputs.

The following deterministic digraph encodings are adopted:

\[
p \rightarrow \AR{بْه}
\]

\[
g \rightarrow \AR{كْ}
\]

\[
v \rightarrow \AR{فْ}
\]

\[
tʃ \rightarrow \AR{تْش}
\]

These sequences are chosen to satisfy:

\[
G(p) \neq G(b)
\]

\[
G(g) \neq G(k)
\]

\[
G(v) \neq G(f)
\]

\[
G(tʃ) \neq G(t), G(ʃ)
\]

Thus injectivity is preserved at the sequence level.

\subsection{Latin-to-Phoneme Parsing}

Let:

\[
\pi : \Sigma_L^* \rightarrow \Sigma_P^*
\]

Parsing rules include maximal matching:

\[
\text{sh} \rightarrow ʃ
\]

\[
\text{ch} \rightarrow tʃ
\]

\[
\text{ng} \rightarrow ŋ
\]

\[
\text{ai} \rightarrow ai
\]

\[
\text{au} \rightarrow au
\]

Singleton mappings apply otherwise.

\subsection{Phoneme-to-Arabic Rewrite Function}

For each \( p \in \Sigma_P \):

\[
p \rightarrow G(p)
\]

Core mappings:

\[
b \rightarrow \AR{ب}
\]

\[
t \rightarrow \AR{ت}
\]

\[
d \rightarrow \AR{د}
\]

\[
k \rightarrow \AR{ك}
\]

\[
f \rightarrow \AR{ف}
\]

\[
ʃ \rightarrow \AR{ش}
\]

\[
m \rightarrow \AR{م}
\]

\[
n \rightarrow \AR{ن}
\]

\[
l \rightarrow \AR{ل}
\]

\[
r \rightarrow \AR{ر}
\]

\[
w \rightarrow \AR{و}
\]

\[
j \rightarrow \AR{ي}
\]

\[
h \rightarrow \AR{ه}
\]

Extended phonemes are represented using the digraph encodings defined above.

\subsection{Injectivity Theorem Under Digraph Encoding}

\textbf{Theorem.}  
Under the digraph encoding strategy, the transliteration mapping

\[
G : \Sigma_P^* \rightarrow (\Sigma_A^{classical})^*
\]

is injective.

\textbf{Proof.}

Assume:

\[
x \neq y
\]

Then there exists position \( i \) such that:

\[
p_i \neq q_i
\]

Since:

\[
G(p_i) \neq G(q_i)
\]

either as single grapheme or digraph sequence, concatenation preserves positional distinction.

Thus:

\[
G(x) \neq G(y)
\]

\hfill \( \square \)

\subsection{Reverse Mapping}

Define inverse mapping:

\[
G^{-1} : (\Sigma_A^{classical})^* \rightarrow \Sigma_P^*
\]

Parsing proceeds according to a precedence rule in which multi-character sequences are identified before single-character mappings are applied. Digraph detection therefore constitutes the first stage of decoding. For example, the sequence \AR{بْه} is interpreted as representing the phoneme /p/, \AR{كْ} is mapped to /g/, \AR{فْ} corresponds to /v/, and \AR{تْش} is resolved as the affricate /tʃ/.

Only after these multi-grapheme combinations have been recognized and resolved does the system apply single-character mappings to the remaining symbols. This ordered procedure prevents premature segmentation and ensures that composite phonological units are not misread as independent consonantal elements.

Maximal-sequence matching ensures deterministic parsing.

\subsection{Algorithmic Implementation}

The system may be implemented as deterministic finite-state transducer:

\[
M = (Q, \Sigma_L, \Sigma_A^{classical}, \delta, q_0, F)
\]

with priority rules for digraph detection.

Time complexity remains:

\[
O(n)
\]

for input length \( n \).

\subsection{Correctness Condition}

Under full vowel marking and digraph preservation:

\[
T_L(T_A(w)) = w
\]

for all phonemically spelled Latin Pidgin input \( w \).

\subsection{Structural Trade-Off}

The removal of Persian letters increases average grapheme length:

\[
|G(p)| \ge 2
\]

for certain phonemes.

However, it preserves:

\[
\text{Injectivity}
\]
\[
\text{Finite-state computability}
\]
\[
\text{Reversibility}
\]

while maintaining strict adherence to Classical Arabic grapheme inventory.

\subsection{Conclusion}

Even under the constraint of Classical Arabic letters only, the transliteration system remains:

\[
\text{Deterministic}
\]
\[
\text{Injective}
\]
\[
\text{Computationally tractable}
\]

The cost is increased graphemic length for non-Arabic phonemes, but structural integrity is preserved.

\section{Appendix K: Arabic-Script Passage and Reverse Transliteration}

\subsection{Arabic Phonetic English}

\begin{quote}
\AR{كَابِتَالِسْتْ يُتُوبِيَانِيزْمْ: ذَا رِبْرِشَنْ أُفْ لَاكْ أَنْدْ إِتْسْ كَلْتْشَرَلْ سِمْبْتُمْزْ}

\AR{إِنْ ذِسْ بَارْتْ أُفْ هِيلِنْ رُولِنْزْ سَايْكُوسِينِمَا، شِي دِلْفْزْ إِنْتُو كَابِتَالِسْتْ آيدِيُولُوجِي أَزْ إِتْ إِنْتِرْسِكْتْسْ وِذْ ذَا سَايْكُوأَنَالِيتِكْ كُونْسِبْتْ أُفْ ``لَاكْ''.}

\AR{شِي فْرَيْمْزْ كَابِتَالِيزْمْ أَزْ أَ سِسْتَمْ بِلْتْ أُونْ دِنَايَلْ أُفْ ذِسْ لَاكْ، رِزَلْتِنْ إِنْ سَايْكُولُوجِيكَلْ أَنْدْ سُوشَلْ كُنْسِكْوِنْسِزْ.}
\end{quote}

\subsection{Direct Phonetic Transliteration}

Capitalist Utopianism: The Repression of Lack and Its Cultural Symptoms

In this part of Heilin Rulins Psychocinema, she delves into capitalist ideology as it intersects with the psychoanalytic concept of “lack.”

She frames capitalism as a system built on denial of this lack, resulting in psychological and social consequences.

\subsection{Standardized English Rendering}

Capitalist Utopianism: The Repression of Lack and Its Cultural Symptoms

In this section of Helen Rollins’ \emph{Psychocinema}, she explores capitalist ideology as it intersects with the psychoanalytic concept of “lack,” understood as an internal gap within the subject.

She argues that capitalism is structured around denial of this constitutive lack, producing both psychological and social consequences.

\subsection{Layered Observations}

The Arabic-script version encodes phonetic English rather than etymological English. Reverse transliteration first yields a phonetic baseline. A secondary smoothing stage restores conventional spelling and standard literary form.

The three layers may be schematized as:

\[
\text{Arabic grapheme} \rightarrow
\text{phonetic English} \rightarrow
\text{standardized English}.
\]

This layered structure confirms that the transliteration system is acoustically grounded while still allowing conventional normalization when required.

\section{Appendix L: Core Nigerian Pidgin Lexicon in Arabic and Latin Scripts}

\begin{longtable}{A p{3.2cm} A p{3.2cm}}
\multicolumn{1}{p{3.2cm}}{\textbf{Pidgin (Arabic)}} &
\textbf{Pidgin (Latin)} &
\multicolumn{1}{p{3.2cm}}{\textbf{English (Arabic)}} &
\textbf{English (Latin)} \\ \hline
\endfirsthead

\multicolumn{1}{p{3.2cm}}{\textbf{Pidgin (Arabic)}} &
\textbf{Pidgin (Latin)} &
\multicolumn{1}{p{3.2cm}}{\textbf{English (Arabic)}} &
\textbf{English (Latin)} \\ \hline
\endhead

هَوْ فَرْ ناو؟ & How far now? & هَوْ آرْ يُو؟ & How are you? \\
آيْ دي & I dey & آيْ أَمْ هِيَرْ & I am here \\
آيْ نو نو & I no know & آيْ دُونْتْ نُو & I don't know \\
وِيتِنْ دِي آبُن؟ & Wetin dey happen? & وَتْ إِزْ جُوِينْ آنْ؟ & What is going on? \\
نا واو & Na wa & ذَتْسْ كْرَيْزِي & That’s crazy \\
آيْ بِلِيفْ سِي & I believe say & آيْ ثِنْكْ ذَتْ & I think that \\
يُو دي كَمْ؟ & You dey come? & آرْ يُو كَمِينْجْ؟ & Are you coming? \\
نُو واهلا & No wahala & نُو بْرُوبْلِمْ & No problem \\
وِايْ ناو؟ & Why now? & وِايْ آرْ يُو دُوِينْ ذِسْ؟ & Why are you doing this? \\
آيْ دي كَمْ & I dey come & آيْ آمْ آنْ مَايْ وِايْ & I’m on my way \\
نُو مينْ & No mean & آيْ دِدْنْتْ مِينْ إِتْ & I didn’t mean it \\
وِايْ يُو دُو آمْ؟ & Why you do am? & وِايْ دِدْ يُو دُو إِتْ؟ & Why did you do it? \\
شِي نو وِيلْ & She no well & شِي إِزْ نَاتْ فِيلِينْجْ وِيلْ & She is not feeling well \\
دِمْ دُونْ كَمْ & Dem don come & ذَيْ آرْ هِيَرْ & They are here \\
بُوِي دِي كَمْ & Boy dey come & ذَ بُوِي إِزْ كَمِينْجْ & The boy is coming \\
آيْ دِي هُنْجَرْ & I dey hunger & آيْ أَمْ هَنْجْرِي & I’m hungry \\
يُو سَابِي بوكْ؟ & You sabi book? & دُو يُو نُو أَنِيثِينْجْ؟ & Do you know anything? \\
وِيتِنْ بِي دِسْ؟ & Wetin be this? & وَتْ إِزْ ذِسْ؟ & What is this? \\
آيْ نو فِيتْ & I no fit & آيْ كَانْ نَاتْ & I can’t \\
آيْ دِي تَايَرْ & I dey tire & آيْ أَمْ تَايَرْدْ & I am tired \\

أَبِجْ & Abeg & بْلِيزْ & Please \\
وَاها لا دِي & Wahala dey & ذِيْرْ إِزْ بْرُوبْلِمْ & There is a problem \\
نُو وَاها لا دِي & No wahala dey & نُو بْرُوبْلِمْ & No problem \\
يُو دُونْ شُورْ & You don sure & آرْ يُو شُورْ؟ & Are you sure? \\
آيْ نو جِيتْ & I no get & آيْ دُونْتْ هَافْ & I don’t have \\
دِمْ دَي بْلَي & Dem dey play & ذَيْ أَرْ نَاتْ سِيرِيُوسْ & They are not serious \\
مَاكْ وِي جُو & Make we go & لِتْسْ جُو & Let’s go \\
كَالْمْ داوْنْ & Calm down & رِلَاكْسْ & Relax \\
يُو سَابِي آمْ؟ & You sabi am? & دُو يُو نُو إِتْ؟ & Do you know it? \\
آيْ نَوْ دِي & I no dey & آيْ أَمْ أَوْكَيْ & I’m okay \\
نَا تْرُو & Na true & ذَتْ إِزْ تُرُو & That is true \\
نَا لايْ & Na lie & ذَتْ إِزْ أَ لايْ & That is a lie \\
يُو دُونْ تْرَاي & You don try & يُو دِدْ وِلْ & You did well \\
آيْ نَوْ فِيتْ دُو آمْ & I no fit do am & آيْ كَانْتْ دُو إِتْ & I can’t do it \\
دِسْ تِنْ سْوِيتْ & This thing sweet & ذِسْ إِزْ جُودْ & This is good \\
آيْ دِي وَاكَا & I dey waka & آيْ أَمْ وُوكِنْجْ & I am walking \\
يُو دِي شَوْ & You dey show & يُو آرْ بْرَاجِنْجْ & You are bragging \\
نَا سُوْ سُوْ & Na so so & أَلْوَيْزْ / تُو مُوشْ & Always / too much \\
آيْ دُونْ تَايَرْ & I don tire & آيْ أَمْ فِيدْ أَبْ & I am fed up \\
مَاكْ آيْ تَلْ يُو & Make I tell you & لِتْ مِي تِلْ يُو & Let me tell you \\

نَا سُوْ آيْ سِي آمْ & Na so I see am & ذَتْسْ هَوْ آيْ سَوْ إِتْ & That’s how I saw it \\
يُو نَوْ هِيِرْ & You no hear & دِدْ يُو نُوتْ هِيَرْ؟ & Did you not hear? \\
آيْ دِي فِينْدْ & I dey find & آيْ أَمْ لُوكِنْجْ فُورْ & I am looking for \\
نَا هُوْ بِي دَتْ؟ & Na who be that? & هُوْ إِزْ ذَتْ؟ & Who is that? \\
نَا وِيتِنْ بِي دَتْ؟ & Na wetin be that? & وَتْ إِزْ ذَتْ؟ & What is that? \\
آيْ نُو سَابِي & I no sabi & آيْ دُونْتْ أَنْدَرْسْتَانْدْ & I don’t understand \\
نَا سُوْ إِهْ بِي & Na so e be & ذَتْسْ هَوْ إِتْ إِزْ & That’s how it is \\
آيْ دِي هِيِرْ يُو & I dey hear you & آيْ أَنْدَرْسْتَانْدْ يُو & I understand you \\
يُو فِيتْ دُو آمْ؟ & You fit do am? & كَانْ يُو دُو إِتْ؟ & Can you do it? \\
آيْ فِيتْ تْرَاي & I fit try & آيْ كَانْ تْرَاي & I can try \\
نَا سُوْ لايْفْ بِي & Na so life be & ذَتْسْ لايْفْ & That’s life \\
آيْ دِي رَنْ & I dey run & آيْ أَمْ إِنْ أَ هُورِي & I am in a hurry \\
يُو دُونْ شُوبْ؟ & You don chop? & هَافْ يُو إِيتَنْ؟ & Have you eaten? \\
آيْ نَوْ وَاكَا رِيتْ & I no waka reach & آيْ دِدْنْتْ جِتْ ذِيْرْ & I didn’t get there \\
دِسْ مَاتَرْ هَارْدْ & This matter hard & ذِسْ إِزْ دِفِيكَلْتْ & This is difficult \\
نَا سَمْتِينْ دِي رُونْ & Na something dey wrong & سَمْثِينْجْ إِزْ رُونْجْ & Something is wrong \\
آيْ نَوْ فِيتْ كِلْ مِيسِلْفْ & I no fit kill myself & آيْ كَانْتْ أُوفَرْوُورْكْ & I can’t overwork \\
يُو دِي يَانْ؟ & You dey yarn? & آرْ يُو تُوكِنْجْ؟ & Are you talking? \\
آيْ دُونْ يَانْ آمْ & I don yarn am & آيْ أَلْرِيدِي سِيدْ إِتْ & I already said it \\
مَاكْ وِي سِي & Make we see & لِتْسْ سِي & Let’s see \\

يُو نُو بِيْتَرْ & You know better & يُو نُو بِيْتَرْ & You know better \\
آيْ دِي تِنْكْ آمْ & I dey think am & آيْ أَمْ ثِينْكِنْجْ أَبَاوْتْ إِتْ & I am thinking about it \\
نَا سَمْتِينْ بِي دِسْ & Na something be this & ذِسْ إِزْ سَمْتِينْجْ & This is something \\
يُو نَوْ دُو آمْ & You no do am & يُو دِدْنْتْ دُو إِتْ & You didn’t do it \\
آيْ دُونْ فُرْجِتْ & I don forget & آيْ هَافْ فُرْجُوتِنْ & I have forgotten \\
دِمْ نَوْ دِي هِيِرْ & Dem no dey hear & ذَيْ أَرْ نَاتْ لِيسِنِنْجْ & They are not listening \\
نَا وِيتِنْ آيْ تَلْ يُو & Na wetin I tell you & ذَتْسْ وَتْ آيْ تُولْدْ يُو & That’s what I told you \\
آيْ دِي كُولْ & I dey cool & آيْ أَمْ كَالْمْ & I am calm \\
يُو دِي فُورْسْ آمْ & You dey force am & يُو أَرْ فُورْسِنْجْ إِتْ & You are forcing it \\
نَا نَاوْ آيْ نُو & Na now I know & نَاوْ آيْ أَنْدَرْسْتَانْدْ & Now I understand \\
آيْ دِي وِيتْ & I dey wait & آيْ أَمْ وِيتِنْجْ & I am waiting \\
يُو نَوْ جِتْ سِنْسْ & You no get sense & يُو أَرْ نَاتْ ثِينْكِنْجْ & You are not thinking \\
دِسْ تِنْ تَايَرْ مِي & This thing tire me & ذِسْ مِيدْ مِي تَايَرْدْ & This made me tired \\
نَا مَايْ فَالْتْ & Na my fault & إِتْ إِزْ مَايْ فَالْتْ & It is my fault \\
آيْ نَوْ دِي فِيلْ آمْ & I no dey feel am & آيْ دُونْتْ لَايْكْ إِتْ & I don’t like it \\
يُو دِي لُوكْ مِي & You dey look me & يُو أَرْ لُوكِنْجْ أَتْ مِي & You are looking at me \\
نَا سُوْ آيْ دِي سِي آمْ & Na so I dey see am & ذَتْسْ هَوْ آيْ سِي إِتْ & That’s how I see it \\
آيْ دِي تِيكْ آمْ & I dey take am & آيْ وِلْ دُو إِتْ & I will do it \\
دِمْ دِي كُولْ مِي & Dem dey call me & ذَيْ أَرْ كَالِنْجْ مِي & They are calling me \\
مَاكْ وِي رِيسْتْ & Make we rest & لِتْسْ رِيسْتْ & Let’s rest \\

يُو دِي سْبِيكْ بِيْجِنْ؟ & You dey speak pidgin? & دُو يُو سْبِيكْ بِيْجِنْ؟ & Do you speak pidgin? \\
وِيتِنْ لَنْجْوِيجْ يُو دِي سْبِيكْ؟ & Wetin language you dey speak? & وَتْ لَنْجْوِيجْ دُو يُو سْبِيكْ؟ & What language do you speak? \\
يُو دِي سْبِيكْ يَوْرُوبَا؟ & You dey speak Yoruba? & دُو يُو سْبِيكْ يُورُوبَا؟ & Do you speak Yoruba? \\
يُو دِي سْبِيكْ إِجْبُو؟ & You dey speak Igbo? & دُو يُو سْبِيكْ إِجْبُو؟ & Do you speak Igbo? \\
يُو دِي سْبِيكْ هَاوْسَا؟ & You dey speak Hausa? & دُو يُو سْبِيكْ هَاوْسَا؟ & Do you speak Hausa? \\
يُو إِنْجْلِيشْ فَايْنْ & Your English fine & يُو إِنْجْلِيشْ إِزْ جْرِيتْ & Your English is great \\
يُو دِي يَانْ وِلْ & You dey yarn well & يُو سْبِيكْ وِلْ & You speak well \\
آيْ لَايْكْ هَاوْ يُو دِي تَاكْ & I like how you dey talk & آيْ لَايْكْ هَاوْ يُو سْبِيكْ & I like how you speak \\
آيْ نِيدْ سُمْبُودِي تُو شُو مِي & I need somebody to show me & آيْ نِيدْ سَمْوَنْ تُو شُو مِي & I need someone to show me \\
شُو مِي هَاوْ تُو مَاكْ فُوفُو & Show me how to make fufu & شُو مِي هَاوْ تُو مَاكْ فُوفُو & Show me how to make fufu \\
شُو مِي هَاوْ تُو مَاكْ إِجُوسِي & Show me how to make egusi & شُو مِي هَاوْ تُو مَاكْ إِجُوسِي & Show me how to make egusi \\
آيْ نَوْ نُو هَاوْ تُو كُوكْ آمْ & I no know how to cook am & آيْ دُونْتْ نُو هَاوْ تُو كُوكْ إِتْ & I don’t know how to cook it \\
دِسْ فُودْ سْوِيتْ وِلْ وِلْ & This food sweet well well & ذِسْ فُودْ إِزْ فِيرِي جُودْ & This food is very good \\
آيْ دِي لِيرْنْ & I dey learn & آيْ أَمْ لِيرْنِنْجْ & I am learning \\
آيْ دِي تْرَاي & I dey try & آيْ أَمْ تْرَايِنْجْ & I am trying \\
تِيكْ آمْ سْمُولْ سْمُولْ & Take am small small & تِيكْ إِتْ سْلُولِي & Take it slowly \\
نَا سُوْ دِمْ تِيكْ تُيچْ مِي & Na so dem take teach me & ذَتْسْ هَوْ ذَيْ تُوتْ مِي & That’s how they taught me \\
آيْ دِي إِنْجُويْ دِسْ & I dey enjoy this & آيْ أَمْ إِنْجُويِنْجْ & I am enjoying this \\
نَا بِيْتَرْ تِنْ بِي دِسْ & Na better thing be this & ذِسْ إِزْ أَ جُودْ تِنْجْ & This is a good thing \\
تَنْكْسْ فُورْ هِلْبْ مِي & Thanks for help me & ثَانْكْ يُو فُورْ هِلْبِنْجْ مِي & Thank you for helping me \\

\end{longtable}

\section{Appendix M: Formal Specification of the Reverse Transliteration and English Normalization System}

\subsection{Overview}

The transliteration framework is organized as a three-tiered representational cascade,

[
\mathcal{A} \rightarrow \mathcal{E}*{\text{phon}} \rightarrow \mathcal{E}*{\text{std}},
]

in which each stage performs a distinct but formally related transformation. The set (\mathcal{A}) designates the Arabic-script phonetic encoding, functioning as the primary acoustic register. The intermediate layer, (\mathcal{E}*{\text{phon}}), represents a phonetic English transcription that preserves pronunciation while adopting a Latin character inventory. The terminal layer, (\mathcal{E}*{\text{std}}), corresponds to standardized orthographic English, where conventional spelling rules are reinstated.

This architecture is formally compositional. The overall transliteration operator may be expressed as

[
T = N \circ R,
]

where

[
R : \mathcal{A} \rightarrow \mathcal{E}_{\text{phon}}
]

denotes the reverse transliteration mapping, and

[
N : \mathcal{E}*{\text{phon}} \rightarrow \mathcal{E}*{\text{std}}
]

denotes the normalization mapping. The composite transformation thus proceeds first through phonetic recovery and subsequently through orthographic regularization, preserving conceptual separation between acoustic representation and literary convention.

\subsection{Stage One: Reverse Transliteration}

Let (\Sigma_A) denote the Arabic grapheme inventory employed in the system, and let (\Sigma_L) denote the Latin phonetic inventory used for English approximation. The reverse transliteration stage is defined by a deterministic reading function

[
R : \Sigma_A^\ast \rightarrow \Sigma_L^\ast,
]

which maps finite sequences of Arabic graphemes to finite sequences of Latin phonetic symbols.

The function (R) operates through a rule-governed decoding procedure. Individual graphemes are mapped to their phonological correlates via grapheme-to-phoneme correspondence rules. Where multi-character clusters represent single phonetic units, digraph resolution is applied, as in the mapping of ت followed by ش to the affricate realized orthographically as *ch*. Vowel quality is determined through explicit diacritic marking, ensuring that vocalic interpretation remains constrained rather than inferred. The resulting phonetic units are then concatenated sequentially to yield a linear Latin transcription.

Provided that vowel diacritics are fully specified and that digraph recognition rules are consistently applied, the mapping (R) is operationally deterministic. Under these conditions, each well-formed input string in (\Sigma_A^\ast) yields a unique phonetic output in (\Sigma_L^\ast), thereby securing reversibility at the level of sound representation.

\subsection{Stage Two: Orthographic Normalization}

The phonetic English layer does not guarantee standard spelling. Therefore define:

\[
N : \Sigma_{L}^\ast \rightarrow \Sigma_{E}^\ast
\]

where $\Sigma_{E}$ is the standardized English orthographic system.

The normalization operator performs a secondary transformation that operates beyond phonetic reconstruction. Whereas the reverse transliteration function preserves acoustic form, normalization introduces conformity to established orthographic and morphological conventions. This stage may involve lexical correction, as when a phonetic rendering such as *utopianizm* is regularized to *utopianism* in accordance with standard spelling. It also includes morphological restoration, reinstating conventional suffixes, possessive markers, or other inflectional features that may not be fully specified at the phonetic layer.

In addition, the operator enforces conventional capitalization patterns and smooths punctuation to align with standard written English usage. These adjustments situate the output within recognized literary norms rather than leaving it at the level of raw phonetic transcription.

Crucially, the operator ( N ) is not purely phonemic in character. Its application presupposes reference to an English lexicon and to morphological rule systems. Unlike the deterministic grapheme-to-phoneme mapping of the reverse transliteration stage, normalization draws upon lexical knowledge and orthographic convention. It therefore functions as a rule-governed but lexically informed interpretive layer, completing the transition from sound-based representation to standardized literary form.

\subsection{Compositional Structure}

The full recovery mapping is:

\[
T = N(R(\cdot))
\]

Thus:

\[
\mathcal{A} \xrightarrow{R} \mathcal{E}_{\text{phon}} \xrightarrow{N} \mathcal{E}_{\text{std}}
\]

The first mapping is phonological; the second is orthographic and lexical.

\subsection{Determinism and Ambiguity}

The mapping $R$ is deterministic under the following conditions:

\[
\text{Full vowel marking} \land \text{Digraph priority parsing}
\]

The mapping $N$ may admit multiple outputs if phonetic spellings correspond to multiple lexical candidates. In practice, contextual constraints restrict ambiguity.

\subsection{Loss and Recovery}

Let:

\[
L = R^{-1}
\]

be forward encoding into Arabic script.

Then:

\[
R(L(x)) = x
\]

holds when $x$ is phonemically spelled.

However:

\[
N(R(L(x))) = x
\]

holds only if $x$ conforms to normalized English spelling conventions.

Thus the system preserves phonetic structure strictly and orthographic convention conditionally.

\subsection{Structural Interpretation}

The transliteration architecture is best understood as a stratified representational system rather than a single-step orthographic substitution. At its foundational level lies acoustic encoding, in which Arabic graphemes function as vehicles for phonological capture. In this first layer, the script is deployed to register sound patterns directly, privileging phonetic approximation over historical spelling conventions or etymological transparency.

The second layer consists of phonetic reconstruction into an English-aligned representation. Here, the encoded sounds are rendered in a form that approximates English phonotactics, but without yet fully conforming to standardized orthographic norms. This intermediate stage serves as a bridge between raw acoustic transcription and conventional spelling, preserving the primacy of pronunciation while moving toward broader legibility.

The third layer introduces standardized literary English. At this stage, conventional orthography is reimposed, aligning forms with established spelling practices. Crucially, this layer is reconstructive rather than foundational. The Arabic-script text does not encode English etymology or standardized spelling at the primary level; it encodes sound. Conventional English spelling therefore emerges as a secondary interpretive stratum, superimposed upon an underlying phonological representation.

\subsection{Conclusion}

The proposed system rests on a set of structural commitments that together define its theoretical coherence. It is phonologically grounded in the sense that its primary representational layer encodes sound rather than historical spelling or etymological lineage. This grounding establishes a direct mapping between acoustic form and graphic sign, minimizing interpretive distortion at the point of transcription.

At the same time, the architecture is algorithmically compositional. Its mappings operate through rule-governed correspondences that can be formalized, iterated, and implemented computationally. The structure is therefore not merely descriptive but generative, capable of systematic application across lexical items without ad hoc adjustment at each instance.

The system is reversible at the phonetic level because its initial encoding preserves sufficient information to reconstruct pronunciation with high fidelity. This reversibility ensures that transliteration does not entail irreversible loss of phonological structure. Yet it remains lexically normalizable at the literary level, allowing forms to be regularized into standardized English orthography where appropriate. Such normalization operates as a distinct interpretive stage rather than as an intrinsic property of the base representation.

The explicit separation between ( R ) and ( N ) secures conceptual clarity within the framework. ( R ) governs acoustic transcription, while ( N ) governs orthographic standardization. By maintaining this distinction, the model avoids conflating phonetic representation with literary convention, thereby preserving analytical precision across representational layers.

\section{Appendix N: Finite-State Formalization of the Transliteration System}

\subsection{Overview}

The transliteration architecture may be modeled as a cascade of deterministic finite-state transducers (FSTs):

\[
\mathcal{A} \xrightarrow{R} \mathcal{E}_{\text{phon}} \xrightarrow{N} \mathcal{E}_{\text{std}}
\]

Each mapping is computable in linear time with respect to input length.

\subsection{Alphabet Definitions}

Let

\[
\Sigma_A = \{\text{Arabic graphemes used in the system}\},
\]

\[
\Sigma_P = \{\text{phoneme symbols for English/Pidgin}\},
\]

\[
\Sigma_E = \{\text{standard English orthographic symbols}\}.
\]

All alphabets are finite sets. In particular, $\Sigma_A$ is restricted to the classical inventory (e.g., \AR{ا}, \AR{ب}, \AR{ت}, \AR{ج}, \AR{ك}, \AR{و}, \AR{ي}), together with diacritics when pedagogical precision is enabled.

\subsection{Transducer 1: Reverse Transliteration}

Define a finite-state transducer

\[
R = (Q_R, \Sigma_A, \Sigma_P, \delta_R, q_0, F_R),
\]

where $Q_R$ is a finite set of states, $\delta_R$ is the transition function, $q_0$ is the initial state, and $F_R \subseteq Q_R$ is the set of accepting states.

The transition function $\delta_R$ implements a deterministic grapheme-to-phoneme mapping composed of four ordered mechanisms. First, single-letter grapheme-to-phoneme correspondence is applied. Second, priority digraph detection is performed for sequences such as \AR{تْش}, which are interpreted as /tʃ/. Third, vowel diacritics (e.g., \AR{َ}, \AR{ِ}, \AR{ُ}) are resolved into their corresponding phonemic values when present. Finally, phoneme outputs are concatenated sequentially to form the resulting string in $\Sigma_P^*$.

Digraph priority is implemented via a bounded lookahead encoded in the state space. For example, upon reading \AR{ت}, the machine transitions to a dedicated intermediate state:

\[
\delta_R(q_0, \AR{ت}) = q_{T}.
\]

If the next input symbol is \AR{ش}, the machine consumes both symbols and emits the affricate:

\[
\delta_R(q_{T}, \AR{ش}) = q_0 \quad \text{with output } /tʃ/.
\]

If the subsequent symbol is not \AR{ش}, then the machine emits /t/ and reprocesses the current symbol in the initial state:

\[
\delta_R(q_{T}, x \neq \AR{ش}) = q_0 
\quad \text{with output } /t/ \text{ and reprocess } x.
\]

This construction guarantees maximal-sequence matching while preserving determinism. No nondeterministic branching is required, and the transduction remains linear in input length.

\subsection{Transducer 2: Normalization Operator}

Define lexical normalization transducer:

\[
N = (Q_N, \Sigma_P, \Sigma_E, \delta_N, q'_0, F_N)
\]

This operator performs a sequence of ordered transformations. It first applies phonetic-to-orthographic replacement, mapping elements of $\Sigma_P$ into conventional forms in $\Sigma_E$. It then applies morphological reconstruction, restoring inflectional patterns and orthographic regularities that may have been neutralized during phonetic normalization. Capitalization rules are subsequently enforced according to standard English conventions, including sentence-initial and proper-noun capitalization. Finally, punctuation smoothing is applied, ensuring that spacing and punctuation conform to conventional typographic norms.

Unlike $R$, $N$ references a finite lexicon:

\[
\mathcal{L} \subset \Sigma_P^\ast
\]

The normalization rule may be written as:

\[
\delta_N(q, w_{\text{phon}}) = (q', w_{\text{std}})
\]

whenever:

\[
w_{\text{phon}} \in \mathcal{L}
\]

If no lexical rule applies, identity mapping is used.

\subsection{Composed System}

The composed transliteration machine is:

\[
T = N \circ R
\]

Since finite-state transducers are closed under composition:

\[
T
\]

is also a finite-state transducer.

Thus the entire Arabic-to-English system is regular and computable.

\subsection{Time Complexity}

For input string length $n$:

\[
R = O(n)
\]

\[
N = O(n)
\]

Therefore:

\[
T = O(n)
\]

The system is linear in input length.

\subsection{Reversibility Conditions}

Let $L$ be the forward encoding operator:

\[
L : \Sigma_P^\ast \rightarrow \Sigma_A^\ast
\]

Then:

\[
R(L(x)) = x
\]

for fully vowel-marked phonemic input.

However:

\[
N(R(L(x))) = x
\]

holds only when $x$ is lexically normalized.

Thus phonemic reversibility is strict; orthographic reversibility is conditional.


\subsection{Structural Significance}

The foregoing formalization establishes that the transliteration system is algorithmically definable in precise automata-theoretic terms. Its operations can be implemented entirely within the class of deterministic finite-state transductions. No context-free machinery is required, since all rewrite rules depend only on bounded local context and finite lookahead. Digraph handling is implemented through explicit state transitions, thereby avoiding nondeterministic branching. The resulting architecture is compositional and modular, in the sense that each stage of processing may be represented as a separate finite-state component whose composition preserves regularity.

\subsection{Conclusion}

The Arabic-script Pidgin and English transliteration framework, as formalized above, constitutes a cascade of deterministic finite-state transducers whose individual components are independently regular and whose composition preserves regularity. Each stage of the pipeline—lexical override, grapheme-to-phoneme conversion, vowel strategy application, orthographic normalization, and reverse mapping—can be expressed as a finite-state relation over finite alphabets. Because the class of regular relations is closed under composition, the global system remains within the same computational class as its parts.

Consequently, the system is

\[
\text{Regular},
\]

\[
\text{Computationally tractable},
\]

and

\[
\text{Formally compositional}.
\]

Regularity entails that all transformations operate with bounded memory and finite lookahead. No recursive stack structures or context-free derivations are required. Computational tractability follows from linear-time processing with respect to input length, ensuring suitability for large-scale digital corpora, real-time input methods, and embedded systems. Formal compositionality guarantees that additional modules—such as morphological smoothing, probabilistic disambiguation layers, or dialect-specific refinements—may be appended without altering the underlying computational classification.

In this respect, the framework aligns with established practices in orthographic engineering, finite-state morphology, and computational phonology. It is compatible with standard finite-state toolchains used in digital linguistic modeling and can be implemented directly in widely adopted transducer architectures. The proposal therefore does not merely describe an orthographic convention; it specifies a formally well-defined computational object situated within the canonical paradigm of modern language technology.

\newpage
\begin{thebibliography}{99}

\bibitem{AbdulazizOsinde1997}
Abdulaziz, Mohamed H., and Ken Osinde (1997).
Sheng and Engsh: Development of Mixed Codes among the Urban Youth in Kenya.
\textit{International Journal of the Sociology of Language} 125, 43--63.

\bibitem{Adegbija2004}
Adegbija, Efurosibina (2004).
The Domestication of English in Nigeria.
In: Edgar W. Schneider et al. (eds.),
\textit{A Handbook of Varieties of English, Volume 1}.
Berlin: Mouton de Gruyter, 20--43.

\bibitem{Aikhenvald2006}
Aikhenvald, Alexandra Y. (2006).
\textit{Serial Verb Constructions in Typological Perspective}.
Oxford: Oxford University Press.

\bibitem{AikhenvaldDixon2006}
Aikhenvald, Alexandra Y., and R. M. W. Dixon (eds.) (2006).
\textit{Serial Verb Constructions: A Cross-Linguistic Typology}.
Oxford: Oxford University Press.

\bibitem{Akinnaso1995}
Akinnaso, F. Niyi (1995).
Literacy, Language and Education in Nigeria.
\textit{Comparative Education Review} 39(1), 68--92.

\bibitem{Ayorinde2013}
Ayorinde, Christine (2013).
Language Policy and the Development of Nigerian Pidgin.
In: Ore Yusuf (ed.),
\textit{Language and Society in Nigeria}.
Ibadan: University Press, 201--220.

\bibitem{Bakker2014}
Bakker, Peter (2014).
Pidgins.
In: Silvia Luraghi and Claudia Parodi (eds.),
\textit{The Bloomsbury Companion to Syntax}.
London: Bloomsbury, 109--130.

\bibitem{BaldiCuly2005}
Baldi, Philip, and Christopher Culy (eds.) (2005).
\textit{Computational Approaches to African Languages}.
Stanford: CSLI Publications.

\bibitem{BeesleyKarttunen2003}
Beesley, Kenneth R., and Lauri Karttunen (2003).
\textit{Finite State Morphology}.
Stanford: CSLI Publications.

\bibitem{Besacier2014}
Besacier, Laurent, et al. (2014).
Automatic Speech Recognition for Under-Resourced African Languages.
\textit{Speech Communication} 56, 85--100.

\bibitem{Bickerton1981}
Bickerton, Derek (1981).
\textit{Roots of Language}.
Ann Arbor: Karoma.

\bibitem{Bickerton1984}
Bickerton, Derek (1984).
The Language Bioprogram Hypothesis.
\textit{Behavioral and Brain Sciences} 7(2), 173--221.

\bibitem{BirdSimons2003}
Bird, Steven, and Gary Simons (2003).
Seven Dimensions of Portability for Language Documentation and Description.
\textit{Language} 79(3), 557--582.

\bibitem{BlanchonLafkioui2017}
Blanchon, Jean A., and Mena Lafkioui (eds.) (2017).
\textit{African Arabic and Contact Linguistics}.
Leiden: Brill.

\bibitem{Blench2006}
Blench, Roger (2006).
\textit{The Languages of Africa}.
Cambridge: Cambridge University Press.

\bibitem{Blench2011}
Blench, Roger (2011).
The Status of Nigerian Pidgin.
\textit{Cambridge Working Papers in Linguistics} 3, 1--27.

\bibitem{Bondarev2017}
Bondarev, Dmitry (2017).
Ajami in West Africa.
In: Fallou Ngom (ed.),
\textit{Muslims Beyond the Arab World: The Odyssey of \textit{Ajam\={i}} and the Mur\={i}diyya}.
Oxford: Oxford University Press, 27--53.

\bibitem{Bybee2010}
Bybee, Joan (2010).
\textit{Language, Usage and Cognition}.
Cambridge: Cambridge University Press.

\bibitem{Cisse2018}
Cissé, Mamadou (2018).
Ajami Writing Systems in West Africa.
\textit{Islamic Africa} 9(2), 135--160.

\bibitem{Comrie1985}
Comrie, Bernard (1985).
\textit{Tense}.
Cambridge: Cambridge University Press.

\bibitem{Coulmas2003}
Coulmas, Florian (2003).
\textit{Writing Systems: An Introduction to Their Linguistic Analysis}.
Cambridge: Cambridge University Press.

\bibitem{Daniels2017}
Daniels, Peter T. (2017).
Script Reform and Orthographic Standardization in Africa.
In: Badreddine Arfi (ed.),
\textit{Writing Systems and Orthography in Africa}.
Dakar: CODESRIA, 45--78.

\bibitem{DanielsBright1996}
Daniels, Peter T., and William Bright (eds.) (1996).
\textit{The World's Writing Systems}.
Oxford: Oxford University Press.

\bibitem{Dwyer1998}
Dwyer, David (1998).
The Use of Arabic Script for African Languages.
In: Karin Barber (ed.),
\textit{Africa's Hidden Histories}.
Bloomington: Indiana University Press, 231--252.

\bibitem{Faraclas1996}
Faraclas, Nicholas (1996).
\textit{Nigerian Pidgin}.
London: Routledge.

\bibitem{Faraclas2013}
Faraclas, Nicholas (2013).
Nigerian Pidgin.
In: Susanne Maria Michaelis et al. (eds.),
\textit{The Survey of Pidgin and Creole Languages}.
Oxford: Oxford University Press, 435--445.

\bibitem{Furniss2004}
Furniss, Graham (2004).
\textit{Literature in Hausa}.
Edinburgh: Edinburgh University Press.

\bibitem{Gambari2015}
Gambari, A. I. (2015).
Ajami Manuscripts and Literacy in Northern Nigeria.
\textit{Sudanic Africa} 26, 87--110.

\bibitem{Hasselblatt2006}
Hasselblatt, Cornelius (2006).
Orthography as a Linguistic Discipline.
\textit{Written Language and Literacy} 9(1), 1--23.

\bibitem{Hasselbring2007}
Hasselbring, Sue (2007).
Hausa Ajami Manuscripts and Orthographic Variation.
\textit{Sudanic Africa} 18, 105--132.

\bibitem{Hockett1958}
Hockett, Charles F. (1958).
\textit{A Course in Modern Linguistics}.
New York: Macmillan.

\bibitem{Holm1988}
Holm, John (1988).
\textit{Pidgins and Creoles, Volume I: Theory and Structure}.
Cambridge: Cambridge University Press.

\bibitem{Holm1989}
Holm, John (1989).
\textit{Pidgins and Creoles, Volume II: Reference Survey}.
Cambridge: Cambridge University Press.

\bibitem{Huber1999}
Huber, Magnus (1999).
\textit{Ghanaian Pidgin English in Its West African Context}.
Amsterdam: John Benjamins.

\bibitem{Huber2008}
Huber, Magnus, and Viveka Velupillai (eds.) (2008).
\textit{Synchronic and Diachronic Perspectives on Contact Languages}.
Amsterdam: John Benjamins.

\bibitem{Hunwick1999}
Hunwick, John (1999).
\textit{Arabic Literature of Africa, Volume 4}.
Leiden: Brill.

\bibitem{Hyman2003}
Hyman, Larry M. (2003).
Segmental Phonology.
In: Derek Nurse and Gérard Philippson (eds.),
\textit{The Bantu Languages}.
London: Routledge, 42--58.

\bibitem{Hyman2007}
Hyman, Larry M. (2007).
Universals of Tone Rules: 30 Years Later.
In: Tomas Riad and Carlos Gussenhoven (eds.),
\textit{Tones and Tunes, Volume 1}.
Berlin: Mouton de Gruyter, 1--34.

\bibitem{KaplanKay1994}
Kaplan, Ronald M., and Martin Kay (1994).
Regular Models of Phonological Rule Systems.
\textit{Computational Linguistics} 20(3), 331--378.

\bibitem{Karan2014}
Karan, Elke (2014).
Orthography Development and Reform.
In: James W. Tollefson and Miguel Pérez-Milans (eds.),
\textit{The Oxford Handbook of Language Policy and Planning}.
Oxford: Oxford University Press.

\bibitem{Karanetal2010}
Karan, Elke, Susan Malone, and John Smalley (2010).
\textit{Developing Orthographies for Unwritten Languages}.
Dallas: SIL International.

\bibitem{Kouwenberg2008}
Kouwenberg, Silvia, and John Victor Singler (eds.) (2008).
\textit{The Handbook of Pidgin and Creole Studies}.
Malden, MA: Wiley-Blackwell.

\bibitem{Labov1972}
Labov, William (1972).
\textit{Sociolinguistic Patterns}.
Philadelphia: University of Pennsylvania Press.

\bibitem{LewisSimonsFennig2016}
Lewis, M. Paul, Gary F. Simons, and Charles D. Fennig (eds.) (2016).
\textit{Ethnologue: Languages of the World}, 19th edn.
Dallas: SIL International.

\bibitem{Lupke2011}
Lüpke, Friederike (2011).
Orthography Development.
In: Peter Austin and Julia Sallabank (eds.),
\textit{The Cambridge Handbook of Endangered Languages}.
Cambridge: Cambridge University Press, 312--336.

\bibitem{McWhorter1998}
McWhorter, John (1998).
Identifying the Creole Prototype.
\textit{Language} 74(4), 788--818.

\bibitem{McWhorter2005}
McWhorter, John (2005).
\textit{Defining Creole}.
Oxford: Oxford University Press.

\bibitem{Michaelis2013}
Michaelis, Susanne Maria, et al. (eds.) (2013).
\textit{The Survey of Pidgin and Creole Languages}, 3 vols.
Oxford: Oxford University Press.

\bibitem{Mufwene2001}
Mufwene, Salikoko S. (2001).
\textit{The Ecology of Language Evolution}.
Cambridge: Cambridge University Press.

\bibitem{Mufwene2008}
Mufwene, Salikoko S. (2008).
\textit{Language Evolution: Contact, Competition, and Change}.
London: Continuum.

\bibitem{Ngom2010}
Ngom, Fallou (2010).
Ajami Scripts in the Senegalese Speech Community.
\textit{Journal of Arabic and Islamic Studies} 10, 1--23.

\bibitem{Ngom2016}
Ngom, Fallou (2016).
\textit{Muslims Beyond the Arab World: The Odyssey of \textit{Ajam\={i}} and the Mur\={i}diyya}.
Oxford: Oxford University Press.

\bibitem{Ngugi1986}
Ngũgĩ wa Thiong'o (1986).
\textit{Decolonising the Mind}.
London: James Currey.

\bibitem{Oyetade2006}
Oyetade, S. O. (2006).
Attitudes to Foreign Languages and Indigenous Languages in Nigeria.
\textit{Sociolinguistic Studies} 1(1), 19--38.

\bibitem{Oyetade2008}
Oyetade, Solomon (2008).
Language Planning in West Africa.
In: Andrew Simpson (ed.),
\textit{Language and National Identity in Africa}.
Oxford: Oxford University Press, 165--187.

\bibitem{Philippson2012}
Philippson, Gérard (2012).
Writing African Languages.
In: Bernd Heine and Derek Nurse (eds.),
\textit{The Oxford Handbook of African Languages}.
Oxford: Oxford University Press, 781--799.

\bibitem{Robinson2006}
Robinson, Clinton (2006).
Language Policy and Literacy in Africa.
\textit{International Journal of Educational Development} 26(4), 379--391.

\bibitem{Rogers2005}
Rogers, Henry (2005).
\textit{Writing Systems: A Linguistic Approach}.
Oxford: Blackwell.

\bibitem{Sanneh1989}
Sanneh, Lamin (1989).
\textit{Translating the Message}.
Maryknoll: Orbis Books.

\bibitem{Schmidt2008}
Schmidt, Annette (2008).
Computational Tools for African Orthographies.
In: Silvia Kouwenberg and John Victor Singler (eds.),
\textit{The Handbook of Pidgin and Creole Studies}.
Malden, MA: Wiley-Blackwell, 570--589.

\bibitem{Schneider2007}
Schneider, Edgar W. (2007).
\textit{Postcolonial English: Varieties Around the World}.
Cambridge: Cambridge University Press.

\bibitem{SIL2004}
SIL International (2004).
\textit{Orthography Development Manual}.
Dallas: SIL International.

\bibitem{Smalley1964}
Smalley, William A. (1964).
Orthography Studies.
\textit{Anthropological Linguistics} 6(2), 1--11.

\bibitem{Smalley1988}
Smalley, William A. (1988).
Orthography as a Social Practice.
\textit{Language in Society} 17(3), 377--393.

\bibitem{Suleiman2013}
Suleiman, Yasir (2013).
\textit{Arabic in the Fray: Language Ideology and Cultural Politics}.
Edinburgh: Edinburgh University Press.

\bibitem{ThomasonKaufman1988}
Thomason, Sarah G., and Terrence Kaufman (1988).
\textit{Language Contact, Creolization, and Genetic Linguistics}.
Berkeley: University of California Press.

\bibitem{Tosco2000}
Tosco, Mauro (2000).
Is Nigerian Pidgin English a Creole?
\textit{Journal of Pidgin and Creole Languages} 15(1), 1--25.

\bibitem{Vydrin2012}
Vydrin, Valentin (2012).
Scripts and Orthographies in West Africa.
\textit{Mandenkan} 48, 3--30.

\bibitem{Winford2003}
Winford, Donald (2003).
\textit{An Introduction to Contact Linguistics}.
Malden, MA: Blackwell.

\bibitem{Winford2008}
Winford, Donald (2008).
Creoles and Contact.
In: Silvia Kouwenberg and John Victor Singler (eds.),
\textit{The Handbook of Pidgin and Creole Studies}.
Malden, MA: Wiley-Blackwell, 3--41.

\bibitem{Zima1974}
Zima, Petr (1974).
\textit{Pidginization and Creolization: The Case of West African Pidgin English}.
Munich: Wilhelm Fink Verlag.

\bibitem{Zwartjes2011}
Zwartjes, Otto (2011).
\textit{Portuguese Missionary Grammars in Asia, Africa and Brazil}.
Amsterdam: John Benjamins.

\end{thebibliography}

\end{document}